% ===========================================================================
%  Qubit & Toffoli Reduction Techniques: A Comprehensive Primer
%  LaTeX source — generated from ecdsafail_optimization_techniques_primer.md
% ===========================================================================
\documentclass[11pt,a4paper]{article}

% ---- Encoding & fonts ----
\usepackage[T1]{fontenc}
\usepackage[utf8]{inputenc}
\usepackage{lmodern}
\usepackage{microtype}

% ---- Page layout ----
\usepackage[top=2.5cm,bottom=2.5cm,left=2.8cm,right=2.8cm]{geometry}
\usepackage{fancyhdr}
\pagestyle{fancy}
\fancyhf{}
\rhead{\small\textit{Qubit \& Toffoli Reduction Techniques}}
\lhead{\small\textit{ecdsa.fail Benchmark Primer}}
\cfoot{\thepage}
\renewcommand{\headrulewidth}{0.4pt}

% ---- Math ----
\usepackage{amsmath,amssymb}

% ---- Colors ----
\usepackage[dvipsnames,svgnames]{xcolor}
\definecolor{codebackground}{gray}{0.95}
\definecolor{codeborder}{gray}{0.75}
\definecolor{tableheadbg}{RGB}{230,240,255}
\definecolor{accent}{RGB}{0,70,140}
\definecolor{warnbg}{RGB}{255,248,230}

% ---- Hyperlinks ----
\usepackage[colorlinks=true,linkcolor=accent,urlcolor=accent,citecolor=accent,
            pdftitle={Qubit \& Toffoli Reduction Techniques: A Comprehensive Primer},
            pdfauthor={ecdsa.fail challenge participants; compiled by Jieyi Long with Claude Code},
            bookmarks=true]{hyperref}

% ---- Code listings ----
\usepackage{listings}
\usepackage{fancyvrb}

\lstset{
  basicstyle=\ttfamily\footnotesize,
  breaklines=true,
  breakatwhitespace=false,
  frame=single,
  rulecolor=\color{codeborder},
  backgroundcolor=\color{codebackground},
  xleftmargin=6pt,
  xrightmargin=6pt,
  framexleftmargin=4pt,
  framexrightmargin=4pt,
  aboveskip=8pt,
  belowskip=6pt,
  showstringspaces=false,
  captionpos=b,
  columns=flexible,
  keepspaces=true,
  tabsize=2,
}

\lstdefinelanguage{Rust}{
  keywords={fn,let,mut,if,else,match,while,for,in,return,pub,use,mod,struct,enum,
             impl,trait,type,where,self,Self,true,false,Some,None,Ok,Err,usize,u8,u16,u32},
  keywordstyle=\color{NavyBlue}\bfseries,
  comment=[l]{//},
  commentstyle=\color{gray}\itshape,
  stringstyle=\color{ForestGreen},
  morestring=[b]",
}

% ---- Tables ----
\usepackage{booktabs}
\usepackage{array}
\usepackage{longtable}
\usepackage{tabularx}
\usepackage{multirow}
\renewcommand{\arraystretch}{1.25}

% ---- Colored boxes ----
\usepackage[most]{tcolorbox}

\tcbuselibrary{breakable,skins}

\tcbset{
  definitionbox/.style={
    breakable,
    enhanced,
    colback=blue!4,
    colframe=accent,
    fonttitle=\bfseries,
    top=4pt, bottom=4pt, left=6pt, right=6pt,
  },
  noteboxstyle/.style={
    breakable,
    enhanced,
    colback=warnbg,
    colframe=orange!60!red,
    fonttitle=\bfseries,
    top=4pt, bottom=4pt, left=6pt, right=6pt,
  },
}

% ---- Enumeration ----
\usepackage{enumitem}
\setlist[itemize]{leftmargin=1.5em,itemsep=1pt,topsep=3pt}
\setlist[enumerate]{leftmargin=1.5em,itemsep=1pt,topsep=3pt}

% ---- Section styling ----
\usepackage{titlesec}
\titleformat{\section}{\large\bfseries\color{accent}}{\thesection.}{0.7em}{}[\color{accent}\titlerule]
\titleformat{\subsection}{\normalsize\bfseries}{\thesubsection.}{0.5em}{}
\titleformat{\subsubsection}{\normalsize\itshape\bfseries}{}{0em}{}
% Number subsections numerically within their section (2.1, 2.2, ... 6.3, ...) —
% the default scheme, restated explicitly to match the in-text references.
\renewcommand{\thesubsection}{\thesection.\arabic{subsection}}

% ---- Utilities ----
\usepackage{xspace}
\usepackage{parskip}
\setlength{\parindent}{0pt}
\setlength{\parskip}{5pt}
\setlength{\emergencystretch}{3em}   % allow extra inter-word stretch to avoid overfull boxes

% ---- Custom commands ----
\newcommand{\ket}[1]{$|#1\rangle$\xspace}
\newcommand{\score}{\textsc{score}\xspace}
\newcommand{\sota}{\textsc{sota}\xspace}
\newcommand{\cls}{\texttt{cls}\xspace}
\newcommand{\pha}{\texttt{pha}\xspace}
\newcommand{\anc}{\texttt{anc}\xspace}
\newcommand{\tof}{Toffoli\xspace}
\newcommand{\avgT}{\texttt{avg\_executed\_Toffoli}\xspace}

\setcounter{tocdepth}{2}

% ===========================================================================
\begin{document}
% ===========================================================================

\begin{titlepage}
\centering
\vspace*{2cm}
{\LARGE\bfseries\color{accent} Qubit \& Toffoli Reduction Techniques\par}
\vspace{0.4cm}
{\Large A Comprehensive Primer\par}
\vspace{1.2cm}
{\large\itshape ecdsa.fail quantum circuit benchmark\par}
\vspace{1.6cm}
{\large\bfseries Contributors\par}
\vspace{0.3cm}
{\normalsize The optimization techniques in this primer were collectively contributed by\par
\textbf{all participants and leaderboard contributors of the ecdsa.fail challenge}.\par}
\vspace{0.5cm}
{\normalsize Summarized and compiled by \textbf{Jieyi Long}, with Claude Code.\par}
\vspace{2.0cm}
\begin{tcolorbox}[definitionbox, title=Audience note]
This primer is written for undergraduate students who have just learned the basics of quantum
computing --- superposition, entanglement, Toffoli/CNOT/Hadamard gates, quantum circuits, and
reversible computation. It explains every optimization technique used in the ecdsa.fail challenge
from circuit-level first principles, with concrete examples from actual research.
\end{tcolorbox}
\vfill
{\small As of 2026-06-27\par}
\end{titlepage}

\tableofcontents
\newpage

% ===========================================================================
\section{Introduction: What Is the ecdsa.fail Challenge?}
\label{sec:intro}
% ===========================================================================

Elliptic-curve cryptography (ECC) secures essentially all of today's digital signatures --- including
the keys that control Bitcoin and Ethereum funds. A large enough quantum computer would break it:
\textbf{Shor's algorithm} can recover a private key from a public key by solving the elliptic-curve
discrete logarithm problem in time polynomial in the key size. The overwhelming majority of that
quantum computation is spent running \textbf{one tiny primitive over and over} --- \emph{elliptic-curve
point addition} on the curve \texttt{secp256k1} (Bitcoin's curve). Make that inner circuit cheaper
and the entire attack gets cheaper.
\href{https://github.com/ecdsafail/ecdsafail-challenge}{\textbf{ecdsa.fail}} is an open, collaborative
competition to build the \textbf{leanest possible reversible quantum circuit for that single point
addition} --- a public, reproducible race to map exactly how expensive (or cheap) breaking ECC really
is.

A submission is a circuit that computes one in-place affine point addition $P \mathrel{+}= Q$ on
secp256k1, and it is scored by its \textbf{spacetime cost: $\texttt{peak\_qubits} \times
\texttt{average\_executed\_Toffoli}$} (lower is better) --- the two resources that dominate cost on a
fault-tolerant quantum computer (\S\ref{sec:score}). Every submission must be \emph{exactly correct
and physically valid}: it is checked against 9{,}024 random test cases for classical-value
correctness, for reversibility (all scratch qubits returned to $|0\rangle$), and for phase
cleanliness, and contributors may only edit the point-addition circuit itself --- the evaluation
harness, validator, and toolchain are locked. The headline goalpost is \textbf{Google Quantum AI's
published resource estimate} for the same task ($\approx$1{,}425 qubits $\times$ $\approx$2.1M Toffoli
$\approx 2.99 \times 10^{9}$ qubit-Toffoli, March 2026). The challenge's starting baseline scored
$\sim$$1.07 \times 10^{10}$; as of late June 2026 the community-driven state of the art sits near
\textbf{1{,}152 qubits $\times$ $\sim$1.36M Toffoli $\approx 1.57 \times 10^{9}$} --- roughly
\textbf{47\% below} Google's number, with the bar still dropping.

What makes the effort unusual is \emph{how} the records are won. Participants --- ranging from
academic researchers to programmers with no quantum-computing background --- increasingly drive the
work with \textbf{AI agent pipelines}: the circuit harness is fed directly into large-language-model
loops that propose, test, and validate optimizations around the clock, preserving their findings as
versioned research notes. The result is a fast-moving catalogue of genuinely clever, circuit-level
tricks for shaving qubits and Toffoli gates. \textbf{This primer is a guided tour of those
techniques}, explained from quantum-circuit first principles for a reader who has just learned the
basics --- covering all three of the challenge's objectives (minimize the spacetime product, minimize
qubits alone, and map the clean (qubit, Toffoli) Pareto frontier; see \S\ref{sec:score}).

% ===========================================================================
\section{What We Are Optimizing --- The Score and Why Toffoli Gates Are Expensive}
\label{sec:score}
% ===========================================================================

\subsection{The score}

\begin{tcolorbox}[definitionbox, title=Score objective (lower is better)]
\[
  \text{score} = \text{peak\_qubits} \;\times\; \text{avg\_executed\_Toffoli}
\]
\end{tcolorbox}

This is the ecdsa.fail benchmark's cost function for a quantum circuit that computes an affine
elliptic-curve point addition on secp256k1 (Bitcoin's curve). The two factors penalize two
different resources:

\begin{itemize}
  \item \textbf{Peak qubits}: the maximum number of qubits simultaneously in use at any instant.
    This is the ``width'' of the circuit, not the total ever allocated. On real hardware, idle
    qubits still occupy physical space on the quantum processor.
  \item \textbf{Average executed \tof}: the Toffoli (CCX) gate count, \emph{weighted by how
    often each gate fires}, averaged over the 9024 test inputs the challenge uses. Not just the
    gate count, but gates~$\times$~firing\_rate.
\end{itemize}

\subsection{Why Toffoli gates specifically?}

On a \textbf{fault-tolerant} quantum computer (the kind needed to break Bitcoin-scale
cryptography), gates fall into two classes:

\begin{itemize}
  \item \textbf{Clifford gates} (CNOT, Hadamard $H$, Phase $S$, Pauli $X/Y/Z$): implemented
    efficiently using the surface error-correcting code at very low physical-qubit overhead.
  \item \textbf{Non-Clifford gates} (\tof / T-gate): require \textbf{magic states} --- specially
    prepared ancilla qubits that take a long time and many physical qubits to distill via
    ``magic state distillation.'' Each \tof gate consumes one distilled magic state.
\end{itemize}

The rough conversion: $1\ \text{Toffoli} = \text{exactly}\ 7\ T\text{-gates}$ (Gosset et al.\
2013, tight lower bound).

So in this regime, \textbf{\tof gates are like currency}: each one costs approximately the same
fixed price (one magic state). The peak qubit count is like ``desk space'' --- idle qubits still
occupy factory floor. The product $\text{peak\_qubits} \times \text{Toffoli}$ captures the
spacetime volume of running the circuit.

\begin{tcolorbox}[noteboxstyle, title=Common metric confusion]
This challenge uses \avgT, NOT:
\begin{itemize}
  \item T-count ($= 7 \times$ Toffoli, a different metric)
  \item T-depth (Selinger 2013 --- a depth trick with 4 ancilla per Toffoli; irrelevant here)
\end{itemize}
The distinction matters: average-executed rewards conditional execution (replay), emitted
T-count does not.
\end{tcolorbox}

\subsection{The current SOTA}

The best known result as of 2026-06-27 (BitWonka, commit \texttt{d44cad3},
trailmix\_ludicrous family):
\[
  1152\ \text{qubits} \times 1{,}364{,}230\ \text{avg Toffoli} = 1{,}571{,}592{,}960\ \text{score}
\]

\subsection{Three objectives, not one}
\label{sec:twocomps}

This is important context for the whole document. The ecdsa.fail challenge is really \textbf{three
optimization tracks}, and a technique that is optimal for one can be irrelevant --- or even
counterproductive --- for another:

\begin{enumerate}
  \item \textbf{Minimize the spacetime volume (the Q$\times$T score competition)} --- the main
    subject of this primer. Minimize the \emph{product}
    $\texttt{peak\_qubits} \times \texttt{avg\_executed\_Toffoli}$. Both resources matter, and every
    move is judged by the exchange rate (\S\ref{sec:exchange}). The SOTA above (1152q $\times$ 1.36M)
    lives in this track.
  \item \textbf{Minimize qubits regardless of Toffoli (the low-qubit competition)} --- a question of
    academic curiosity: \emph{how few qubits can a correct point-add possibly use?} Toffoli count is
    \emph{unconstrained}. This track maps the qubit \textbf{floor}; the deepest technique here is EEA
    register sharing (\S\ref{sec:regshare}), reaching 851q $\to$ 829q (far below anything
    score-competitive). Under this objective, spending 8$\times$ more Toffoli to save 10\% of qubits
    is \emph{free}, because Toffoli never enters the scoring.
  \item \textbf{Explore the (Q, T) Pareto frontier} --- especially the \emph{useful corner} where
    \textbf{both} Q and T are lower than the published resource estimates for a \textbf{single
    secp256k1 point addition}: \textbf{$\leq$1175 qubits / $\sim$2.69M Toffoli} ($= 2^{21.36}$), the
    Babbush et al.\ (Google Quantum AI) space-optimized figure tabulated in Schrottenloher 2026
    (arXiv:2606.02235, Table 1; original arXiv:2603.28846). \emph{(Watch the units: the often-quoted
    $\leq$1200q / $\leq$90M Toffoli is the \textbf{full} ECDLP/Shor run --- $\sim$28 windowed point
    additions plus overhead --- whereas the scored circuit here is \textbf{one} point addition, so
    the per-point-addition figure is the correct reference.)} A point is Pareto-optimal if you cannot
    lower one of Q or T without raising the other. The goal here is not a single score but a
    \emph{clean, forkable basis curve} at each qubit count (see the 1153q$\to$1133q Pareto bases,
    \S\ref{sec:paretopush}). To stay reusable, these circuits \textbf{deliberately disable the most
    overfit lever --- dead-CCX deletion} (\S\ref{sec:density}), whose hard-coded gate-index lists are
    tied to one exact op stream. They are \emph{not} fully value-exact, though: they still use the
    route's calibrated \textbf{approximations} (comparator-width narrowing, carry truncation, and at
    the lowest rungs active-width trimming --- all island-exact, \S\ref{sec:density}), so each still
    needs its own nonce hunt. The point is a \emph{cleaner, more forkable} reference than a
    dead-CCX-stacked SOTA --- not an all-inputs-provable one.
\end{enumerate}

These are genuinely different games. Most of this primer optimizes track 1, but several techniques
(register sharing, the clean Pareto bases) only make sense under tracks 2 and 3. Keep the three
objectives separate in your head: many ``is this worth it?'' questions only have an answer once you
know \emph{which} track you are on.

% ===========================================================================
\section{The Reversibility Constraint --- Why Quantum Circuits Are Hard to Optimize}
\label{sec:reversibility}
% ===========================================================================

\subsection{You cannot delete information}

Quantum circuits must be reversible. This is not a convention --- it is required by quantum
mechanics. Every gate in a quantum circuit must have an inverse (you can always ``run it
backwards''). This means:

\begin{itemize}
  \item \textbf{You cannot overwrite a qubit.} If qubit $A$ holds value $x$ and you want to
    compute $f(x)$ into $A$, you cannot just replace it. You'd lose $x$ forever, breaking
    reversibility.
  \item \textbf{You cannot erase a qubit.} If a qubit is entangled with the rest of the system
    (as intermediate results always are), setting it to \ket{0} forcibly would violate unitarity.
\end{itemize}

The only way to ``free'' a qubit is to \textbf{uncompute} it --- run the computation that created
its value in reverse, returning it exactly to \ket{0}. Once it's \ket{0}, it can be freed.

\subsection{Consequence: intermediate values cost qubits until uncomputed}

Every time your circuit computes an intermediate result (a carry bit, a partial product, a
comparison flag), it must allocate a fresh qubit. That qubit remains \textbf{live} --- occupying
a slot in the peak count --- from the moment it's computed until the moment it's uncomputed.

\textbf{Example}: An $n$-bit addition $a + b$ computes $n$ carry bits, one per bit position. Each
carry qubit is live from when it's computed until the carry chain is reversed. If you compute the
addition at step 100 and uncompute it at step 200, those $n$ carry qubits are live across the
entire steps 100--200, contributing to the peak at any instant in that range.

This is why almost every qubit-reduction technique boils down to:
\textbf{make intermediate values die sooner} (so they aren't live at the peak) \textbf{or don't
materialize them at all}.

\subsection{Reversibility also doubles Toffoli counts}

In a classical circuit, you compute $X$, use $X$, then discard $X$. In a reversible circuit, you
compute $X$ (costs \tof), use $X$, then \emph{uncompute} $X$ (costs the same \tof again). So
every logical AND operation normally costs $2\times$ --- once forward, once backward.

This is the central pain point that Measurement-Based Uncomputation (\S\ref{sec:mbu}) solves.

% ===========================================================================
\section{The Challenge Circuit --- Elliptic Curve Point Addition}
\label{sec:circuit}
% ===========================================================================

\subsection{What it computes}

The circuit computes one \textbf{affine elliptic curve point addition}: given quantum register
$P = (P_x, P_y)$ and classical constant $Q = (Q_x, Q_y)$, compute $P + Q = R$ in place, where
the arithmetic is modulo the secp256k1 prime $p = 2^{256} - 2^{32} - 977$.

The mathematical operations, in order:
\begin{lstlisting}
dx = Px - Qx              (mod p subtraction)
dy = Py - Qy
lambda = dy / dx          (modular inverse + multiply: the expensive GCD step)
new_Px = lambda^2 - Px - Qx
new_Py = lambda*(Px - new_Px) - Py
\end{lstlisting}

\subsection{Why modular inverse is the bottleneck}

Computing $dy/dx \bmod p$ requires computing $dx^{-1} \bmod p$, the modular inverse of $dx$.
There is no simple formula --- it requires running the \textbf{Extended Euclidean Algorithm (EEA)}
or equivalent, which takes $\approx530$ steps of binary GCD (divstep) operations. Each step involves:
\begin{itemize}
  \item A comparison (swap decision: is $|u| > |v|$?)
  \item A conditional swap of two 256-bit registers
  \item A conditional subtraction of two 256-bit registers
\end{itemize}
This accounts for roughly 90\% of the circuit's \tof count.

\subsection{How an in-place modular inverse actually works: the dialog transcript}
\label{sec:dialog}

A classical EEA tracks two ``B\'ezout coefficient'' registers alongside the operands to build the
inverse. Doing that reversibly would carry several extra 256-bit registers through all $\sim$530
steps --- very expensive in qubits. The trick that makes the circuit affordable (Schrottenloher
Algs 2--4, descended from the Khattar--Gidney ``Dialog'') is to \textbf{split the algorithm into a
cheap construction pass and a cheap reconstruction pass that share a small recorded tape}:
\begin{itemize}
  \item \textbf{Pass 1 --- GCD construction (records the tape).} Drive the binary-GCD pair $(u, v)$
    from $(p, dx)$ toward $(1, 0)$. The only data-dependent information at each step is the small
    \textbf{decision symbol} (subtracted?, swapped?, shift) --- a few bits. Record those onto a
    \textbf{transcript tape}. This pass uses \emph{non-modular} arithmetic with $\sim$n free ancilla,
    so it is cheap. It \emph{consumes} $dx$ into the tape.
  \item \textbf{Pass 2 --- B\'ezout reconstruction (replays the tape).} Read the tape \textbf{in
    reverse}, maintaining one \emph{modular} pair $(r, s)$ and applying only linear updates
    \textbf{controlled solely by the recorded bits}: $s = 2s \bmod p$; if subtracted: $s \mathrel{+}= r$;
    if swapped: $\mathrm{swap}(r,s)$. No comparisons, no operand-dependent branching --- just replay.
\end{itemize}

\textbf{The elegant part (why no explicit inverse is ever built).} Because the replay updates are
\emph{linear}, the seed you feed into pass 2 comes straight out. Seeding $(r, s) = (1, 0)$
reconstructs $dx^{-1}$; seeding $(r, s) = (y, 0)$ instead yields $(0,\; y\cdot dx^{-1} \bmod p)$
\textbf{directly} --- the inversion \emph{and} the multiply-by-$y$ happen together, with the inverse
never materialized as its own register. And reading the tape \textbf{forward vs reversed flips
between $y\cdot dx$ and $y\cdot dx^{-1}$} --- the same machinery is a modular multiply or a modular
divide by replay direction. This is how the SOTA does both EC slope steps with \textbf{one inversion
engine and one tape} (\S\ref{sec:trailmix}).

\textbf{Where the peak lives --- and the one knob that sets the Pareto split.} The qubit peak is
\emph{not} during construction (operands shrink, freeing room for the tape); it is during
\textbf{reconstruction}, where the compressed transcript ($\sim$2.355n) and the two n-bit modular
registers $r, s$ co-reside:
\[
  \text{peak} \approx \text{compressed transcript} + 2\,\text{modular registers} + O(\sqrt{n})
  \approx 4.355n.
\]
So \emph{both} score factors are set by the reconstruction step --- that is where transcript
compression (\S7.8) and the vented adders (\S9.1) are aimed. And one clean knob sits behind the whole
(Q, T) Pareto curve: spending \textbf{n extra ancilla during reconstruction} lets you use a cheap
Gidney adder (2n Toffoli) instead of a CDKM adder (3n) --- that single choice \emph{is} the
space-optimized-vs-gate-optimized split in both Babbush and Schrottenloher.

\textbf{Jump-GCD (fewer steps).} A plain binary GCD does one reduction per trailing-zero bit. A
\textbf{jump (Stein-style) divstep} strips up to \texttt{jump=2} trailing zeros at once, so the
algorithm finishes in fewer steps --- directly fewer adder/comparator firings, hence fewer Toffoli.
The cost is a slightly larger symbol alphabet to record (5 symbols instead of 3), which the codec
absorbs. The SOTA runs \texttt{jump=2} (258 steps).

\subsection{The key design choice: Q is classical}

The second point $Q$ is a \textbf{fixed constant} (the secp256k1 generator point, known at circuit
compile time). The current SOTA keeps $Q$ as classical bit-patterns (\texttt{BitId}s) rather than
quantum registers --- \emph{and} materializes it into a transient quantum temp only at off-peak steps
--- eliminating \textbf{512 qubits} from the peak: the single largest qubit reduction in the circuit's
history, developed as a technique in \S\ref{sec:qclassical}.

\subsection{Why secp256k1's prime is cheap: pseudo-Mersenne reduction (the math background)}
\label{sec:pmmath}

Every arithmetic step works \textbf{modulo} $p = 2^{256} - 2^{32} - 977$. Doing $\bmod\,p$ naively
means dividing by a 256-bit number --- a full multiprecision division, expensive in any circuit. The
whole reason this challenge is approachable is that $p$ is a \textbf{pseudo-Mersenne (Solinas) prime}
--- it sits just below a power of two. Write
\[
  p = 2^{256} - c, \qquad c = 2^{32} + 977 \quad\text{(only 33 bits wide)}.
\]

\textbf{A quick word on the names} (they describe a family of ``reduction-friendly'' primes):
\begin{itemize}
  \item \textbf{Mersenne prime} --- $p = 2^k - 1$. Reduction is \emph{trivial}: $2^k \equiv 1$, so you
    fold the high bits straight onto the low bits with one add. (Too rare to land on a secure 256-bit
    curve, but the ideal everything else approximates.)
  \item \textbf{Pseudo-Mersenne prime} --- $p = 2^k - c$ with $c$ \emph{small} (here 33 bits). Now
    $2^k \equiv c$, so the fold is ``high bits $\times c$, then add.'' \emph{Pseudo-Mersenne}
    emphasizes that \textbf{$c$ is small}, which keeps the $H\cdot c$ multiply cheap.
  \item \textbf{Solinas prime} (a.k.a.\ \emph{generalized Mersenne}, Jerome Solinas, 1999) --- $c$ is
    not just small but a \textbf{sparse signed sum of a few powers of two}. For secp256k1,
    $c = 2^{32} + 977 = 2^{32} + 2^{10} - 2^{5} - 2^{4} + 1$ (see NAF, \S9.7) --- only $\sim$5 signed
    terms. \emph{Solinas} emphasizes that \textbf{$c$ is sparse}, so the $H\cdot c$ ``multiply'' is
    really a handful of \emph{shifted adds/subtracts}, not a general multiplication.
\end{itemize}
secp256k1's prime is \textbf{both at once}, and the two properties cut different costs: being
pseudo-Mersenne (small $c$) means \emph{one} fold suffices; being Solinas (sparse $c$) means that fold
is a few shift-and-adds. Throughout this primer ``the Solinas fold'' and ``pseudo-Mersenne
reduction'' refer to the same operation --- folding the top bits down by $+c$ --- and ``Solinas
constant'' / ``the fold constant'' both mean $c$ (called $f$ in the code).

The one identity everything rests on is
\[
  2^{256} \equiv c \pmod{p}
\]
(because $2^{256} - c = p \equiv 0$). In words: \textbf{anything at bit position 256 or higher is
congruent, mod $p$, to the same thing multiplied by the tiny constant $c$ and dropped back into the
low bits.} So a wide value $H\cdot 2^{256} + L$ reduces as
\[
  H\cdot 2^{256} + L \;\equiv\; L + H\cdot c \pmod{p}.
\]
You never divide by $p$; you \textbf{``fold'' the overflow $H$ back down} by multiplying it by the
33-bit $c$ and adding --- a couple of short operations instead of a division. This single fact is the
source of the entire secp256k1-vs-generic-prime cost gap: in Schrottenloher's accounting the
constant-adder is \textbf{15\% of all Toffoli for secp256k1 but 34\% for a generic prime}, and
modular doubling \textbf{8\% vs 24\%} (paper Table 3; secp space-opt $2^{21.19}$ vs generic
$2^{21.78}$).

The intuition to carry forward: a modular reduction is normally ``subtract $p$ until in range,'' but
for a pseudo-Mersenne prime it becomes \textbf{``fold the top bits into the bottom bits,''} cheap
precisely because $c$ is small. \S\ref{sec:pmarith} turns this into the actual reversible gadgets
(mod-double, mod-add, the $+f$ fold); \S9.7 (NAF) makes the fold constant itself cheaper.

% ===========================================================================
\section{Measuring Progress --- The Peak Qubit Profile and Hot Spots}
\label{sec:measuring}
% ===========================================================================

\subsection{The live-qubit-vs-time profile}

To reduce the peak, you first need to know \textbf{where} the peak occurs. Plot how many qubits
are live at each moment in the circuit timeline:

\begin{lstlisting}
qubits  |
live    |         ##########
        |    ####################
        |  #########################
        |################################
        +--+----+--+------+--+----+----> time
                    ^
                    peak = 1152 here (the "binder")
\end{lstlisting}

The highest point is the \textbf{peak}, and the phase executing at that height is the
\textbf{binder}. Only operations alive at the binder instant matter for qubit reduction.

\textbf{Key tool}: \texttt{TRACE\_PEAK=1 TRACE\_PHASE\_ACTIVE=1} in the build flags. This
instruments every \texttt{alloc\_qubit}/\texttt{free\_qubit} call and prints the name of the
phase executing when the maximum is hit.

\subsection{The plateau problem}

After shaving the single tallest binder, you often discover the peak hasn't moved --- because
several \emph{independent} phases all reach the same height. This is a \textbf{co-binder plateau}:

\begin{lstlisting}
qubits  |
live    |    #####   #####     ####
        |  #########################
        |  #########################  <- plateau at 1152 (Phase A, B, C all tie here)
        +----+--+----+---+--+---+---> time
\end{lstlisting}

The reported peak doesn't drop until \textbf{every} plateau member is individually pushed below the
current height. This changes the strategy fundamentally: you must work in ``plateau mode'' ---
find a local qubit reduction for each co-binder, then combine them all.

\subsection{Emitted vs average-executed Toffoli}

\begin{itemize}
  \item \textbf{Emitted}: the raw count of CCX gates in the circuit, regardless of whether they fire.
  \item \textbf{Average-executed}: each CCX gate's contribution weighted by $\Pr[\text{control fires}]$,
    averaged over all 9024 test inputs. A CCX that only fires 5\% of the time contributes only 0.05.
\end{itemize}

This distinction is critical because \textbf{conditional replay} (\S\ref{sec:replay}) can
dramatically reduce average-executed without changing emitted at all.

% ===========================================================================
\section{Qubit Reduction: The Core Loop}
\label{sec:qubitloop}
% ===========================================================================

Before diving into individual techniques, understand the general loop:

\begin{enumerate}
  \item \textbf{Measure} --- get the live-qubit-vs-time trace; find the peak height and the binder phase.
  \item \textbf{Inventory} --- list every qubit alive at the binder instant.
  \item \textbf{Classify} --- label each qubit by category (table below).
  \item \textbf{Brainstorm} --- for each non-essential qubit at the binder, propose a shave.
  \item \textbf{Cost} --- compute the \tof overhead vs qubits saved. Compare to the break-even exchange rate (\S\ref{sec:exchange}).
  \item \textbf{Re-measure} --- apply the cheapest safe shave. The peak moves to a new binder. Repeat.
\end{enumerate}

\subsection{Qubit classification at the binder}

\begin{center}
\begin{tabular}{p{3.2cm}p{5.5cm}p{3.5cm}}
\toprule
\textbf{Class} & \textbf{What it holds} & \textbf{Shrinkable?} \\
\midrule
Essential data & Inputs/outputs the binder reads or writes & Rarely \\
Scratch & Temporary arithmetic results (carries, partial products) & Often \\
Transcript / log & GCD branch decisions recorded for reversibility & Often --- compress the log \\
Passenger & Allocated before the peak, consumed after, \emph{idle at binder} & Yes --- relocate its endpoints \\
Alias-candidate & Restorable duplicate of a value already live elsewhere & Yes --- borrow as dirty scratch \\
Truncatable & Bits provably zero or irrelevant to the final output & Yes \\
\bottomrule
\end{tabular}
\end{center}

\subsection{Binary vs graded failure modes}

A critical distinction before applying any technique:

\begin{itemize}
  \item \textbf{Graded-failure levers} (carry truncation, comparator narrowing): correctness
    degrades \emph{continuously} with aggressiveness. A small fraction of inputs go wrong. You can
    iterate on the parameter and search for a nonce where all 9024 test inputs avoid the failures.

  \item \textbf{Binary-failure levers} (register live-range holes, transcript compression,
    structural reuse): correctness is all-or-nothing. Either the circuit is value-exact on all
    inputs, or it's wrong on essentially \emph{all} inputs. There is no graded near-miss ladder
    to climb. You must get the implementation mathematically correct before searching for nonces.
\end{itemize}

\begin{tcolorbox}[noteboxstyle, title=The trap]
Applying the ``tweak-and-triage'' loop to a binary-failure lever.
A structural width cut that's not value-exact will produce all-dirty garbage on every input, not
near-misses.
\end{tcolorbox}

% ===========================================================================
\section{Qubit Reduction Techniques (1--20)}
\label{sec:qubittech}
% ===========================================================================

\subsection{Live-Range Holes: Uncompute Early, Recompute Later}
\label{sec:liverange}

\textbf{The headline move.} If a qubit holds a value $V$ that the peak instant does NOT need (it
was computed before the peak and will be consumed after), free it before the peak and rebuild it
after.

\begin{lstlisting}
Timeline view:
Before:  ...[V computed]=====[V held live across peak]=====[V consumed]...
After:   ...[V computed][uncompute V]  (peak happens)  [recompute V][V consumed]...
                                        ^^^^^^^^^^^
                                  V is absent at peak
\end{lstlisting}

\textbf{Mechanics}: Emit the reverse of the gates that built $V$ (returning its qubits to \ket{0}
and freeing them), run the peak-owning phase with that freed headroom, then re-emit the forward
gates to restore $V$ for its later consumer.

\textbf{Cost}: For a value that cost $C_0$ Toffoli to compute:
\begin{itemize}
  \item Save $k$ qubits at the peak
  \item Spend $2C_0$ extra Toffoli (one uncompute + one recompute)
  \item The trade wins if $2C_0 < k \times (\text{avgT} / \text{peak\_q})$ (the break-even rate)
\end{itemize}

\textbf{Real examples}:
\begin{itemize}
  \item \textbf{1211$\to$1193}: A GCD-apply scratch block was uncomputed during a shift sub-phase
    that didn't use it, then reacquired after. Saved 18 qubits.
  \item \textbf{1166$\to$1165} (\texttt{TLM\_PARK\_ODD\_U0}): \texttt{u[0]} is provably 1 (GCD
    odd-u invariant). Apply $X \to$ \ket{0} to zero it, \texttt{loan\_zero\_qubit} to return the
    lane, then \texttt{restore\_known\_one} (apply $X$ again) to restore. Near-zero \tof cost,
    $-1$ qubit.
\end{itemize}

\textbf{The keep-alive dual}: Sometimes the uncompute+recompute round-trip creates a wider
transient intermediate than just holding the value. In that case, \emph{keep} the value live
across the middle operation. Choose whichever is narrower: holding footprint vs recompute peak
transient.

\subsection{Passenger Relocation}

A \textbf{passenger} is a qubit that was allocated early and consumed late, but does nothing at
the peak --- it just sits idle.

\textbf{The cheaper fix}: Don't uncompute and recompute. Just \textbf{allocate later} (right
before first use) and \textbf{free earlier} (right after last use). If both endpoints fall
outside the peak window, the passenger vanishes from the binder count at zero \tof cost.

\textbf{Check passengers first} before resorting to live-range holes --- relocation is free.

\subsection{Range-Bound a Register to Drop a Guard Bit}

Some registers need an extra ``guard bit'' (sign or overflow indicator) only because an
intermediate could transiently go negative or overflow. If you can prove the value always stays
in a non-negative range by reordering operations, delete the guard bit entirely.

\textbf{Example} (benhuang Lever B): The fold quotient register had bit 33 (a sign guard) because
Solinas reduction terms could make the partial sum dip below zero. Reorder to apply positives
first. Now the sum is always in $[0, 2^{33})$. Allocate 33 bits instead of 34.
\textbf{Net: $-1$ qubit at zero \tof cost.} This produced the 1221$\to$1220 drop.

\textbf{General rule}: In $+A - B$ sequences, applying positives first and subtracting last
keeps partial sums non-negative, eliminating sign-guard bits.

\subsection{Truncation: Keep Only the Bits That Matter}

Carry chains, high words of partial sums, and intermediate registers often hold more bits than
actually affect the final result.

\textbf{Per-step scheduling}: In a 530-step loop (the GCD), a single constant truncation width
must be set safe for the worst iteration. A per-step schedule $\text{width}(\text{step})$ lets
each iteration be exactly as tight as it can tolerate --- saving qubits only at the steps that
are actually at the peak.

\textbf{Two truncation flavors}:
\begin{itemize}
  \item \textbf{Structural truncation (density-neutral)}: prove from circuit analysis that the
    dropped bits are always zero on all inputs.
  \item \textbf{Empirical truncation (island-exact)}: verify the dropped bits are never 1 on the
    9024 challenge inputs specifically. Requires re-hunting a nonce after each change.
\end{itemize}

\subsection{Free-and-Recompute a Cheap Value (The 1153$\to$1152 Breakthrough)}

When a dead-qubit scan finds \textbf{no} unused qubits at the peak, don't conclude
``the floor is structural.'' Ask instead:

\begin{quote}
  \textit{Is any held value a cheap function of the other values already live at the peak?}
\end{quote}

If yes, \textbf{free} that qubit, let the peak phase use its slot, then \textbf{recompute} the
value from the live qubits afterward for negligible cost.

\textbf{The cy0 trick} (\texttt{d44cad3}, BitWonka): The qubit \texttt{cy0} holds the carry-in
at bit 0. No dead qubits found by scan. But because the secp256k1 prime constant $f[0] = 1$:
\[
  \texttt{cy0} = \texttt{control\_bit}\ \text{AND}\ (\text{NOT}\ \texttt{final\_a0})
\]
Both \texttt{control\_bit} and \texttt{final\_a0} are already live at the peak. So:
\begin{enumerate}
  \item Uncompute \texttt{cy0} using 2 CNOT-equivalent gates ($\approx 0$ Toffoli)
  \item Return the freed lane via \texttt{loan\_zero\_qubit}
  \item After the peak phase: recompute \texttt{cy0} from \texttt{control\_bit} and
    \texttt{final\_a0} (2 gates)
\end{enumerate}
Applied 40 times (once per binding fold iteration): $40 \times 4$ CNOTs $= 0$ Toffoli.
Saves 1 qubit. The 1153$\to$1152 drop.

\textbf{Generalized lesson}: After a dead-qubit scan fails, scan again for ``whose value is a
CHEAP FUNCTION of other live qubits?'' The second question catches what the first misses.

\subsection{Dynamic-Width Registers}

In a GCD algorithm, operands naturally shrink as computation proceeds. After $k$ steps of binary
GCD, the operands $u$ and $v$ have lost at most $k$ bits from their high ends. A naive circuit
allocates full 256-bit registers for all 530 steps.

\textbf{Dynamic-width registers}: At each step $i$, determine the actual live bit-width and only
allocate that many qubits. Free the high bits as they become known-zero.

\textbf{Ablation results} (from the Shrunken-PZ q980 track):
\begin{lstlisting}
Source-level dynamic PZ (universal schedule):     1027 qubits
+ thin support-tuned schedule (per-step widths):   990 qubits  (-37q)
+ counter/rotation metadata trimming:              983 qubits   (-7q)
+ quotient cap:                                    980 qubits   (-3q)
\end{lstlisting}

\textbf{Thin schedule (island-exact)}: An even tighter version samples actual GCD widths over many
random inputs and builds per-step bounds tight to that distribution. This is island-exact ---
the thin schedule may be too narrow for inputs outside the sampled set. Requires nonce hunting.

\subsection{Known-Constant Teardown Before the Peak}

When the GCD terminates, certain registers are in known, predetermined states: $A = 0$, $B = 1$,
$\texttt{ca} = p$, $q = 0$. These are \textbf{classical constants} --- not quantum information.

\textbf{What to do}: XOR the known constant \emph{out of} the register (zeroing it) before the
peak computation, then XOR it \emph{back in} afterward.

\textbf{Example}: Register \texttt{ca} holds $p = 2^{256} - 2^{32} - 977$ at termination. Before
the lambda multiply: XOR $p$ out of \texttt{ca}, returning it to \ket{0}. Free the register. Run
the lambda multiply with the freed headroom. After the lambda multiply: XOR $p$ back in for
uncomputation.

\textbf{Toffoli cost}: XOR (CNOT) operations are Clifford gates --- free.
The teardown and recreation cost 0 Toffoli. This is a pure qubit win.

\subsection{Transcript / Log Compression}
\label{sec:transcript}

\textbf{The GCD must record its decisions.} At each of the 530 divstep iterations, the algorithm
makes data-dependent decisions: did we swap? did we subtract? These decisions are recorded in a
\textbf{transcript} (also called ``dialog tape'') because they're needed to run the GCD backward
during uncomputation. With one qubit per decision, the raw transcript is $\approx772$ bits long.

\textbf{Compression techniques}:

\textbf{1. Base-5 codec} (trailmix\_ludicrous): Each 3-step window of $(\texttt{subtracted},
\texttt{swap}, s_2)$ has 8 possible bit patterns, but only 5 are reachable. Encode each 3-step
window in $\lceil \log_2(5^3) \rceil = 7$ bits instead of 9 raw bits.
\begin{lstlisting}
Raw:        1 + 3*257 = 772 qubits
Compressed: 1 Step0 (2b) + 85 Triples (7b each) + 1 Pair tail (5b) + t1 flag (1b) = 603 qubits
\end{lstlisting}

\textbf{2. K5 exact 11-bit codec}: Enumerate all reachable 5-step transcript windows
exhaustively. If there are exactly $2^{11} = 2048$ reachable patterns, encode each in 11 bits.
Information-theoretically optimal. Implemented via a borrowed-ancilla 17-CCX reversible codec.
Converts 1203q$\to$1193q.

\textbf{3. Three-step tail codec} (32 symbols): The last few GCD steps are heavily constrained.
Enumerate the exact reachable set (32 symbols), encode in 5 bits. The 1185q$\to$1170q drop.

\textbf{4. Streaming decompression}: Instead of decompressing the entire transcript before use,
decompress one window at a time, use it, then immediately clear that window. Only one window is
``open'' at a time, reducing peak overhead from the codec.

\textbf{The hard constraint}: The codec must be a \textbf{bijection} --- a perfect 1-to-1 map
between reachable patterns and codes. Every codec must include a self-test proving bijectivity over
the full reachable support.

\subsection{Plateau Decomposition: Lower Every Co-Binder Together}

When multiple independent phases all tie at the same peak height, you're in plateau mode.
Lowering one phase alone doesn't change the global peak.

\textbf{The protocol}:
\begin{enumerate}
  \item List ALL phases tied at the current peak (not just the first one). Build a ``co-binder ledger.''
  \item Find a \emph{local} qubit reduction for each binder independently.
  \item Verify each local reduction doesn't accidentally raise a neighboring phase.
  \item Combine all compatible local reductions in one package.
  \item Re-measure the combined package.
\end{enumerate}

\textbf{Real example --- the q1170 plateau} (9-way co-bind, all tied):
\begin{center}
\begin{tabular}{p{4.5cm}p{5.5cm}p{3cm}}
\toprule
\textbf{Phase} & \textbf{Local fix} & \textbf{Effect} \\
\midrule
Solinas square & smaller square segment notch & square peak 1170$\to$1169 \\
apply final ripple & low-qubit final ripple variant & ripple 1170$\to$1162 \\
apply double/reverse halve & fold carry parking + per-step map & 1170$\to$1169 \\
special overflow/underflow folds & fold parking + per-step map & 1170$\to$1169 \\
non-final apply ripple & source-as-carry suffix for one lane & 1170$\to$1169 \\
\bottomrule
\end{tabular}
\end{center}
Each was independently developed, then combined. Only when all binders were locally reduced
did the global peak drop to 1169.

\subsection{Dynamic Headroom Clamp}

Define a \textbf{circuit-wide qubit ceiling} and have every arithmetic primitive clamp its
carry/chunk width to stay under it:

\begin{lstlisting}[language=Rust]
fn target_qubit_headroom(circ: &Circuit) -> usize {
    TLM_TARGET_Q - circ.active_qubits()
}

fn varchunk_adder(circ: &mut Circuit, bits: usize) -> ... {
    let k = min(scheduled_k, target_qubit_headroom(circ));
    // ... use k as the carry chunk width ...
}
\end{lstlisting}

\textbf{Effect}: No single call can exceed \texttt{TLM\_TARGET\_Q} total qubits.
\textbf{Lower \texttt{TLM\_TARGET\_Q} by 1 $\to$ peak drops by 1}, as long as the narrower
adder is still correct.

\textbf{The tradeoff}: Tighter ceilings force narrower carries $\to$ more carry-chunk splits $\to$
more Toffoli. The break-even exchange rate governs when to tighten.

\textbf{Crucial nuance --- the clamp is only a \emph{ceiling}, not a lane-freeing lever.} Lowering
\texttt{TLM\_TARGET\_Q} declares an aspiration; it does nothing unless there are \emph{actually
freeable lanes} for the narrower arithmetic to fit into --- ``a clamp with no freed lanes just
panics or fails to fit.'' So a clamp drop must be \textbf{paired with companion levers that genuinely
eliminate live qubits} at the peak. The \textbf{1156 $\to$ 1153} drop is the textbook example --- it
stacked three coordinated levers:

\begin{enumerate}
  \item \textbf{The clamp} itself: \texttt{TLM\_TARGET\_Q} 1156 $\to$ 1152 (the binding peak lands
    at 1153, one above the aspiration --- the irreducible co-resident set).
  \item \textbf{Zero-cin fold} (\texttt{TLM\_FOLD\_CHUNK\_ZERO\_CIN}): the fold-chunk loop's
    persistent carry-in ancilla \texttt{cin0} is \emph{no longer allocated at all} --- the first
    chunk's carry-in is provably 0, so a new ``zero-cin boundary-erase'' path
    (\lstinline{fold_boundary_erase_zero_direct}) does the carry cleanup without it. This is the
    actual freed lane (a live-range elimination, cf.\ \S\ref{sec:lazycarry}) that the clamp reclaims.
  \item \textbf{FFG vent-count cap} (\texttt{TLM\_FFG\_MAX\_G = 47}): caps the FFG half-product
    fold's clean-vent count $g$ so fewer transient vent qubits co-reside at the peak
    ($\texttt{capped\_g} = \min(\texttt{scheduled\_g}, \texttt{cap})$).
\end{enumerate}

\textbf{And the qubit drop alone is not enough to \emph{win}.} Pushing the peak 1156 $\to$ 1153
makes the clamped adders run on tighter chunks $\to$ more carry work $\to$ avgT went \emph{up} by
$+11{,}369$ (to 1{,}392{,}603), i.e.\ $\sim$3{,}790 Toffoli/qubit, far above break-even. That
``clean q1153 base'' was \textbf{rejected} until Toffoli levers (the cinc cascade replacement
\S\ref{sec:cinc} + iterated dead-CCX \S\ref{sec:deadccx}) were stacked back on top to carry it over
the line --- a perfect illustration of the product-race rule (\S\ref{sec:exchange}): a qubit drop is
only real on the score once enough Toffoli is recovered to pay for it.

\subsection{Lazy Carry Liveness}
\label{sec:lazycarry}

A specific instance of live-range shortening for carry-in qubits. The carry-in qubit was being
allocated before a chunk loop and held live across all $N$ iterations, even though it was only
needed at the loop boundaries.

\textbf{Fix} (\texttt{TLM\_FOLD\_CHUNK\_LAZY\_CIN0}):
\begin{lstlisting}
Before: carry_in allocated -> live across all N chunk iterations -> used only at boundaries
After:  carry_in allocated inside boundary_erase() -> used -> freed inside boundary_erase()
\end{lstlisting}
The carry-in qubit's liveness is reduced to just the boundary operation.
Peak saved: 1 qubit per chunk at 0 Toffoli cost.

\subsection{In-Place Decode: Hold Blocks Compressed}

For persistent register blocks that live across the peak, encode them compressed and decode
\emph{in place} only at the single use site.

\textbf{Before}: Full-width raw K2-pair block (6 lanes) held live across the whole peak.
\textbf{After}: 5 compressed cells + one zero lane stored at rest; decode in place only at the
use site (allocate the 6 lanes transiently, use them, re-encode to 5 cells). During the rest of
the circuit, only 5 cells are allocated. This is \texttt{DIALOG\_GCD\_K2\_APPLY\_INPLACE\_RAW\_BLOCK=1},
the 1218$\to$1215 qubit drop.

\subsection{Borrowed-Carry Fusion}
\label{sec:borrowcarry}

Standard carry/borrow vectors allocate fresh qubits. But if an adjacent scratch register is
currently known \ket{0}, borrow its qubits as the carry vector instead:

\begin{lstlisting}[language=Rust]
borrowed_const_fold_carries(circ, need, borrowed_carries);
// Uses borrowed[..need] as carry vector; no fresh allocation
// At exit: borrowed restored to |0>
\end{lstlisting}

\textbf{Net effect}: No fresh qubits allocated for the carry prefix. The peak doesn't rise.

\textbf{Hard requirement}: The borrowed qubits must be genuinely \ket{0} at the borrow point and
restored to \ket{0} before the borrower reads them again as data. Getting this wrong causes
silent data corruption. (\texttt{DIALOG\_GCD\_SPECIAL\_FOLD\_BORROW\_CARRIES=1}.)

\subsection{Passenger Ghosting via HMR}

\textbf{Standard passenger relocation}: Free before peak; recompute after. Requires the value
to be recomputable from live data (costs Toffoli).

\textbf{HMR ghost} (more aggressive): Measure the passenger qubit in the Hadamard ($X$-basis).
The qubit is freed. The measurement result (a classical bit $m$) is recorded. Later, reconstruct
the value algebraically and apply a classically-conditioned $Z$ correction to fix the phase.

\textbf{Why this works}: An $X$-basis measurement of $|\psi\rangle$ projects it into $|+\rangle$
or $|-\rangle$. The measurement outcome $m$ tells you which. If $m = 1$ (got $|-\rangle$), a
phase $(-1)$ was applied to the computational-basis components --- the CZ correction cancels this
exactly. After the correction, the measured qubit is factored out and can be freed without
information loss.

\textbf{In the q980 EC point addition}:
\begin{itemize}
  \item $dy$ (256-bit) computed before the GCD inversion. After GCD, $\lambda$ and $dx$ are live.
  \item The identity $dy = \lambda \cdot dx$ holds by construction.
  \item HMR-ghost $dy$ before the GCD peak (measure it, record $m$).
  \item At the end: reconstruct $dy = \lambda \cdot dx$, apply phase correction via $m$.
  \item Freed 256 qubits across the peak at the cost of one modular multiply for reconstruction.
\end{itemize}

\subsection{EEA Register Sharing --- The Low-Qubit Frontier}
\label{sec:regshare}

\textbf{The deepest qubit-reduction idea, and the headline technique for the low-qubit competition}
(\S\ref{sec:twocomps}, track 2 --- minimize peak qubits, Toffoli unconstrained). This is what drives
the lowest qubit counts ever recorded on the challenge: the 948q $\to$ 851q $\to$ 829q descent.
Because that competition does not score Toffoli at all, register sharing is free to spend gates
lavishly in exchange for lanes --- exactly the trade it makes.

\textbf{The naive EEA register budget.} The extended Euclidean algorithm carries four full-width
256-bit registers through all $\sim$530 steps:
\begin{itemize}
  \item $A$, $B$ --- the two operands (the shrinking $u$, $v$)
  \item $CA$, $CB$ --- their two B\'ezout cofactors (the running $a$, $b$ with $a\cdot dx + b\cdot p = \gcd$)
\end{itemize}
Naively that is $4 \times 256 = 1024$ qubits just for the inversion core, before any scratch.

\textbf{The sharing observation} (Proos--Zalka 2003, the origin of ``inversion dominates point-add
space''): the operands and their cofactors are \emph{complementary in size}. As the algorithm runs,
operand $A$ loses high bits (it shrinks toward 1) while its cofactor $CA$ grows from zero. At every
step the sum $\mathrm{bitlen}(A) + \mathrm{bitlen}(CA)$ is bounded by roughly one word. So $A$ and
$CA$ can be \textbf{packed into a single physical 256-bit register}: $A$ in the high part, $CA$ in
the low part, with a small \emph{bit-length tracker} recording where the boundary sits. As $A$
shrinks and $CA$ grows, the boundary simply slides --- but the total never exceeds one word.

\begin{lstlisting}
Naive:   [====== A ======][== CA (mostly zero) ==]   two separate 256-bit words
         [====== B ======][== CB (mostly zero) ==]

Shared:  [== A ==|== CA ==]   ONE word; boundary slides right as A shrinks, CA grows
         [== B ==|== CB ==]   ONE word; a 9-bit length tracker remembers the split
\end{lstlisting}

Four full registers collapse to \textbf{two ``work words'' + a handful of small length/control
trackers}. This is the structural source of Luo et al.\ 2026's compact exact reversible inversion at
$3n + 4\lfloor \log_2 n \rfloor$ qubits (1333q for $n=256$ as a standalone inversion). In the
Shrunken-PZ track, packing $A\Vert CA$ and $B\Vert CB$ plus dirty-catalytic borrowing and
gate-hosting drove the inversion core down to
\lstinline{REGISTER_SHARED_PAPER_INVERSION_QUBITS = 3*256 + 52 = 820} qubits, giving an
\textbf{851-qubit peak} for the whole point-add --- a 97-qubit drop below the prior 948q record, and
the 851$\to$829q race then squeezed out another 22.

\textbf{Where the Toffoli goes (and why that's fine here).} Packing is not free in gates: a shared
word has a \emph{data-dependent boundary} --- ``where does $A$ end and $CA$ begin?'' depends on
values computed at runtime. So every arithmetic step must first \textbf{locate and realign that
boundary with a full-width comparator}, and because the circuit is reversible, that comparator is
recomputed \emph{before and after} each hosted gate:

\begin{lstlisting}
Cost to realize ONE useful Toffoli on a packed word:
   toggle_gate ; ccx ; toggle_gate
   where toggle_gate ~= a full-width comparator ~= 12 x width ~= 3,075 Toffoli at width 256
\end{lstlisting}

Across 530 divsteps this lands the 851q route at $\sim$464.5M average Toffoli (vs $\sim$54.8M for
the unshared 948q route). \textbf{Under the low-qubit objective this is a non-issue} --- Toffoli does
not enter the score, so spending $\sim$8$\times$ more of it to buy 97 qubits is precisely the
intended trade. (For contrast, under the \emph{Q$\times$T product} objective the very same packing
would be $\sim$250$\times$ over break-even, so this technique belongs squarely to the low-qubit
track, not the score track. The two objectives have genuinely different optimal constructions --- see
\S\ref{sec:exchange} for the exchange-rate math behind that statement.)

\begin{tcolorbox}[noteboxstyle, title={How to read the 851q figure}]
851q is a carefully-scoped \textbf{lower-bound witness} --- a research milestone that says ``this
width looks reachable,'' with its modeling assumptions documented openly in the source. It is exactly
the kind of result objective 2 (\S\ref{sec:twocomps}) is meant to surface. Three things are worth
knowing so you read the number for what it is:
\begin{itemize}
  \item \textbf{The deepest figure rests on a \emph{classical} feasibility model.}
    \lstinline{register_shared_eea.rs} computes on \texttt{U256}/\texttt{U512} integers --- a
    classical simulation rather than an emitted gate stream. Its job is to confirm the packing
    invariant $l_t + 1 + l_q + \mathrm{bitlen}(r) \le n+3$ holds at all 1{,}479 steps --- i.e.\ that
    operand and cofactor genuinely \emph{fit in one word}, establishing the packing is feasible in
    principle. The authors flag this transparently: \emph{``an analysis oracle $\ldots$ passing this
    oracle does not establish a challenge-valid qubit count''} (i.e.\ a feasibility model, to be
    paired with a hardened emission before it counts as a finalized circuit).
  \item \textbf{The lowest counts use gate-hosting whose zero-claims are \emph{sampled} rather than
    proven} (see \S\ref{sec:gatehost}). The route is honest about which guarantees are still open:
    \lstinline{q945_local_hosts.rs} notes the reachable-state zero claims are not yet certified and
    lists the remaining ``Q946 hardening'' prerequisites; \lstinline{q944_gate_host_lifecycle.rs}
    verifies the hosting exhaustively \emph{at small widths} and extrapolates to 256 bits.
  \item \textbf{The status is labelled, not hidden.}
    \texttt{Q949\_ROBUST\_ENVELOPE\_FRESH\_VALIDITY\_CLAIMED = false} and the \texttt{proof\_status}
    fields mark these as lower-bound explorations rather than finalized circuits --- good hygiene for
    a frontier-mapping result.
\end{itemize}
In short: the $0/0/0$ run check confirms the \textbf{outputs are correct on the hunted island}, and
the 851/829q figures show \textbf{how low the width can plausibly go}. The lowest count backed by a
fully-emitted, island-validated circuit with \emph{no} outstanding hosting assumptions is the
\textbf{948q} route --- the natural ``validated'' milestone to cite --- while 851q/829q are the
further-reaching witnesses on top of it. Either way, the durable, fully-sound takeaway is the
structural mechanism: operand/cofactor packing.
\end{tcolorbox}

\textbf{The lesson.} Register sharing is \emph{the} route below $\sim$948 qubits, and the
Proos--Zalka / Luo-2026 mechanism is the structural source worth knowing. It is the right tool when
your objective is pure qubit count. The broader point generalizes both ways: \textbf{know which
competition you are playing before pricing a lever} --- a move that is a clear win for the low-qubit
objective (trade unlimited Toffoli for lanes) can be a clear loss for the Q$\times$T objective, and
vice versa.

\subsection{Gate-Hosting: Borrowing an Idle Lane Instead of Allocating}
\label{sec:gatehost}

\textbf{A workhorse peak-reduction technique --- it powers the 948 $\to$ 851 drop
(\S\ref{sec:regshare}), and it doubles as a clean illustration of the value-exact vs island-exact
distinction at the level of a \emph{single borrowed qubit}.}

\textbf{The idea.} Every reversible gadget needs \textbf{scratch ancilla} (a multi-controlled-X
needs a lane to accumulate the AND of its controls; a carry chain needs carry bits; a comparator
needs a borrow bit). Allocating a fresh $|0\rangle$ for each is what drives the peak up. But at the
instant you need scratch, some lane \emph{elsewhere} is often idle --- the high bits of a shrunk
number, a register between its last read and next write, an idle metadata lane.
\textbf{Gate-hosting borrows that idle lane as scratch, then restores it before its owner needs it
again} --- so the ancilla costs no new peak qubit. The scratch is ``hosted'' inside a register that
belongs to something else (the codebase has explicit per-step tables naming which \texttt{A/B/CA/CB/Q}
lane+bit hosts each gate). This generalizes \S\ref{sec:borrowcarry} (borrowed-carry fusion) to
\emph{any} gate's scratch.

\textbf{The hot-desking intuition.} Qubits are desks in a crowded office. Allocating a fresh ancilla
$=$ renting a new desk (the office gets more crowded). Gate-hosting $=$ doing your hour of work at a
colleague's empty desk while they are at lunch, then clearing it spotless before they return --- no
new desk, the peak never rose.

\textbf{Two flavors:}
\begin{itemize}
  \item \textbf{Clean hosting} --- the borrowed lane is known $|0\rangle$. Use it directly, restore
    it to $|0\rangle$.
  \item \textbf{Dirty hosting} --- the borrowed lane holds unknown live data $\psi$. You can still
    borrow it via a trick that disturbs $\psi$ and then exactly un-disturbs it (the Barenco
    ``surrounded'' pattern):
\begin{lstlisting}
ccx(c0, c1, psi)        <- scribble onto the borrowed lane
ccx(psi, c2, target)    <- do the real work using it
ccx(c0, c1, psi)        <- un-scribble: psi restored to its original value
ccx(psi, c2, target)
\end{lstlisting}
    The two extra gates restore $\psi$ perfectly. Cost: 4 Toffoli instead of 3 --- a textbook
    \textbf{$-1$ qubit, $+1$ Toffoli} trade (great when the objective does not score Toffoli).
\end{itemize}

The principled, scaled-up form of dirty hosting is \textbf{conditionally-clean ancillae}
(Khattar--Gidney 2024, ``Laddered Toggle Detection''): a method to borrow dirty bits \emph{as if
they were clean}, with no allocation and no separate toggle-detection pass. It is what lets a
multi-controlled-X run in $\sim$3n Toffoli with only $\log^* n$ truly-clean ancilla, and it underlies
the cheap MCX/\texttt{cinc} primitives (\S\ref{sec:cinc}). When you need a lot of scratch and have
only dirty lanes, this is the formal tool.

\textbf{It is a compile-time bet, not a runtime scan.} A quantum circuit is static --- you cannot
peek at a qubit mid-run to check ``is this $|0\rangle$?'' (measuring would collapse it). So nothing
searches for zeros at runtime. The hosting choices are \textbf{frozen into the gate list at build
time}: the compiler asserts \emph{``at step 437, lane 12 of register B will be $|0\rangle$, so host
the scratch there,''} and the circuit runs that exact plan on every input. The whole question is
whether that compile-time claim was \emph{true}.

\textbf{Proven-zero vs sampled-zero (value-exact vs island-exact, one qubit at a time).} The
zero-claims come in two grades:
\begin{itemize}
  \item \textbf{Proven zero (sound, density-neutral):} the lane is zero for a \emph{structural}
    reason --- e.g.\ the high bits of an operand the GCD has provably shrunk, or a known-constant
    register. Safe for all inputs.
  \item \textbf{Sampled zero (island-exact, risky):} the claim ``this lane is $|0\rangle$ here'' is
    established by an \textbf{empirical census} --- run many sample inputs, observe the lane was
    always zero, \emph{assume} it always is. This is the grade the aggressive 851q borrows rely on.
    ``The desk was empty every time I checked'' is not ``the desk is always empty.''
\end{itemize}

\textbf{How it breaks.} If an input arrives (one not in the sample) that makes a sampled-zero host
lane \emph{non-zero} at borrow time, the hosted gate scribbles onto live data and the circuit
silently returns a wrong answer --- reversibility offers no safety net. The reason the shipped
circuit still scores $0/0/0$ is that the \textbf{nonce/test island is hunted to dodge exactly those
inputs}: the 9024 graded inputs are chosen so none trips a bad borrow. So it never breaks on its own
test set --- because the test set was fitted to the circuit, not because the circuit is universally
correct.

\textbf{Takeaway.} Gate-hosting is a perfectly sound, standard peak-reduction technique \textbf{when
the host lane's state is proven}; it becomes an island-overfit assumption the moment the state is
only \emph{sampled}. The same value-exact vs island-exact split from \S\ref{sec:density} applies
here at the granularity of a single borrowed $|0\rangle$.

\subsection{Reset-Bounded Qubit-ID Compaction (Register Allocation for Qubits)}
\label{sec:compact}

\textbf{A free, value-exact post-pass that closes the gap between the \emph{true} peak and the
\emph{scored} peak.} The benchmark charges \textbf{the largest qubit \emph{id} that appears in the
circuit}, not the true count of simultaneously-live qubits. The builder reuses freed lanes, but if
its numbering is even slightly loose --- a temporary gets a high id when a lower one was already free
--- the \emph{max id} can sit above the real concurrent peak. That slack is pure wasted score.

\textbf{The fix is classic register allocation, applied to qubit ids.} After the whole op stream is
built, run a compaction pass (\lstinline{compact.rs}, introduced in the 1133 Pareto push):
\begin{enumerate}
  \item Cut the timeline at every \textbf{unconditional reset} (\texttt{R}) or measurement
    (\texttt{Hmr}) --- these destroy a lane's value, freeing its id for reuse from that point.
  \item Build the \textbf{lifetime interval} of every non-IO temporary lane (first use $\to$ last use
    within a reset-bounded segment).
  \item \textbf{Interval-color} those intervals onto the fewest ids (a min-heap of free colors keyed
    by lifetime end --- textbook interval-graph coloring), while \textbf{pinning the four IO
    registers} so the evaluator still finds its inputs/outputs.
\end{enumerate}

The result is the same circuit (every gate identical) renumbered so $\texttt{max id}$ equals the
coloring optimum. It is \textbf{value-exact and density-neutral} --- it changes only the \emph{labels}
on wires, never the computation --- so it stacks freely and needs no nonce re-hunt.

\textbf{Why this is the right mental model.} A qubit id is exactly a \emph{register} in a compiler:
lanes never live at the same time can share one id, just as non-overlapping variables share a machine
register. Lifetime-coloring is how compilers minimize register pressure, and it is the provably
optimal way to minimize $\texttt{max id}$ for a fixed gate schedule. Whenever a circuit's reported
qubit count exceeds its measured concurrent peak (check with \texttt{TRACE\_PEAK}), this pass recovers
the difference for free. (In the 1133 push the lever is present but gated off --- 1133 came from
active-width trimming, not compaction --- so it is an untapped drop-in for any future rung whose
\texttt{max id} exceeds its true peak.)

\subsection{Windowed (Chunked) Materialization --- the \texttt{F\_CUT} Dial}
\label{sec:windowed}

\textbf{The biggest peak lever of the 1693$\to$1411 era.} Many operations need a \emph{materialized}
operand --- a full-width quantum register built from something cheaper. The classic case: to add the
quantum-controlled value $f = \texttt{ctrl}\cdot a$ into an accumulator (where $a$ is 256 bits), the
naive circuit first builds all 256 bits of $f$, then runs a 256-bit ripple adder. Two wide things are
now live at once --- the materialized $f$ \emph{and} the adder's carry lane --- and their sum is the
peak.

\textbf{The fix: never materialize the whole operand at once.} Split the add into chunks. For each
chunk: materialize only that slice of $f$, run a small adder over just that slice, then
\textbf{immediately clear the slice with a measurement (HMR)} before moving to the next chunk. Only
one chunk's worth of $f$ is ever live, and the boundary carry between chunks is cleaned up with a
short truncated comparator against the source bits.

\begin{lstlisting}
Naive:   build all 256 bits of f=ctrl.a  ->  256-bit ripple add   (f + carry both wide -> tall peak)
Chunked: per block:  build slice of f  ->  add slice  ->  HMR-clear slice  ->  fix boundary carry
         (only one slice of f live at a time -> much shorter peak)
\end{lstlisting}

\textbf{Why it's a \emph{dial}, not a one-shot.} The chunk boundary (\texttt{F\_CUT},
\texttt{APPLY\_CHUNKED\_F\_BLOCKS}) is tunable: each $+1$ narrows the materialized slice and lowers
the peak by $\sim$1--2 qubits, at the cost of one more boundary comparator (a few Toffoli). So you
``sink'' the apply phase to the next floor by turning the dial, then re-tune it after every
\emph{other} lever lands. This single knob, re-cut at nearly every step from 1572q down to 1411q, did
most of the era's peak work. It is the materialized-operand version of venting and per-step
truncation: \textbf{process a wide value a window at a time so its full width never co-resides at the
peak.}

\subsection{Shift-Free Scaling via Modular Doublings}
\label{sec:shiftfree}

A subtle qubit trap hides in ``multiply by a power of two.'' Computing $x \cdot 2^k \bmod p$ the
obvious way --- $\texttt{shift\_left}(k) \to \texttt{op} \to \texttt{shift\_right}(k)$ --- needs a
\textbf{spill register} to catch the $k$ bits that fall off the top during the shift, plus
overflow/sign flags. In the Solinas $2^{32}$ fold of the square tail, that parked $\sim$24 persistent
flag qubits and pinned the peak.

\textbf{The fix uses a number-theoretic identity:} modulo $p$, multiplying by 2 \emph{is} a modular
doubling (\texttt{mod\_double}, \S\ref{sec:pmarith}) --- a shift that immediately folds its overflow
back via $+f$, leaving nothing to spill. So $x \cdot 2^k \bmod p$ can be done as \textbf{$k$ successive
modular doublings} (and $\div 2^k$ as $k$ modular halvings), value-identical to the shift route but
allocating \textbf{no spill register and no flags}. It trades a handful of extra Toffoli for $\sim$24
freed qubits (\texttt{KARA\_SOL\_SHIFT22\_DOUBLES}, the 1558q drop). The lesson: a ``free'' classical
operation (a bit-shift) can be quietly \emph{expensive} in a reversible circuit because the
shifted-out bits must be stored, not discarded --- and re-expressing it in the modular arithmetic you
already have can dodge that storage entirely.

\subsection{Keep Classical Constants Off the Peak (the $-512$q Q-classical lever)}
\label{sec:qclassical}

\textbf{The single largest qubit reduction in the circuit's history ($-512$q), and a technique worth
stating on its own.} The second EC point $Q = (Q_x, Q_y)$ is a \emph{fixed, compile-time-known
constant} (the generator). Naively you would load its two coordinates into 256-bit quantum registers
and let them ride through the point-add. But two 256-bit quantum registers that sit \textbf{live
across the GCD peak} cost \textbf{512 qubits at the binder} --- the most expensive place to spend them.

\textbf{The lever has two halves, and you need both:}
\begin{enumerate}
  \item \textbf{Keep the constant \emph{classical}.} Store $Q_x, Q_y$ as classical bit-patterns
    (\texttt{BitId}s), not qubits. Arithmetic against a classical operand is cheaper anyway --- adding
    a known constant is a pattern of CNOT-like gates with no quantum operand (\S9.9, the Toffoli side
    of the same fact).
  \item \textbf{Materialize it \emph{off-peak} only.} A few steps genuinely need $Q$ as a transient
    quantum value (the coordinate folds). Load it into a short-lived quantum temp \textbf{only at those
    off-peak steps}, use it, and free it immediately --- so the 512 qubits exist only during
    low-occupancy phases and \textbf{never co-reside with the GCD peak}.
\end{enumerate}

Half (1) alone is not enough: earlier circuits already knew $Q$ was a constant, yet still materialized
it as resident quantum operands across the peak --- which is exactly why the dialog-GCD frontier
stalled at the 1168q wall (\S\ref{sec:history}). It was the \emph{ludicrous} restructuring (an
$a$-independent affine $\lambda$-law that arranges the whole add as two GCD passes, so $Q$ is only ever
needed transiently at off-peak coordinate steps) that finally delivered half (2). tob-joe's note calls
this \textbf{``the decisive product-min lever ($-512$q).''}

\textbf{The general principle:} \emph{a value that is (a) a compile-time constant and (b) only needed
at off-peak moments should be kept classical and reconstructed into a transient quantum temp at its
point of use --- never held resident at the peak.} The savings scale with the constant's width and how
long it would otherwise be live; for the two 256-bit $Q$ coordinates across the full GCD, that is the
whole 512.

% ===========================================================================
\section{Toffoli Reduction: The Core Loop}
\label{sec:tofloop}
% ===========================================================================

\begin{enumerate}
  \item \textbf{Measure} --- get the \tof count; find the \textbf{hot-spot}. Note emitted vs average-executed.
  \item \textbf{Classify} --- identify what \emph{kind} of Toffoli these are (table below).
  \item \textbf{Pick the matching move} --- the kind determines the technique.
  \item \textbf{Cost it} --- depth-neutral? qubit cost? worth it at the exchange rate?
  \item \textbf{Re-measure} --- the hot-spot migrates; repeat.
\end{enumerate}

\textbf{Classifying Toffoli by kind:}
\begin{center}
\begin{tabular}{p{6.5cm}p{6cm}}
\toprule
\textbf{Kind of Toffoli at the hot-spot} & \textbf{The matching move} \\
\midrule
Carry-uncompute (reverse carry chain returning scratch to \ket{0}) & Vent via MBU (\S\ref{sec:mbu}) \\
Self-uncomputing scratch (comparators built then cleanly reversed) & Conditional replay (\S\ref{sec:replay}) \\
Data-controlled arithmetic (\texttt{cadd}, fold-sub on live data) & Fusion / avoid materialization \\
Majority gadgets (3-CCX majority in carry logic) & Ancilla-free majority (\S\ref{sec:majority}) \\
Redundant final cleanup & Skip it (\S\ref{sec:skipclean}) \\
Full intermediate products / wide rows & Avoid materialization (\S\ref{sec:witness}) \\
Adjacent algebraic stages (add/negate/subtract chains) & Algebraic fusion (\S\ref{sec:fusion}) \\
\bottomrule
\end{tabular}
\end{center}

% ===========================================================================
\section{Toffoli Reduction Techniques (1--19)}
\label{sec:toftech}
% ===========================================================================

\subsection{Measurement-Based Uncomputation (MBU)}
\label{sec:mbu}

\textbf{The \tof lifecycle problem}: An AND gate costs 4 $T$-gates (1 Toffoli) to compute. The
uncomputation normally costs another 4 $T$-gates. Computing + uncomputing one AND $= 2$ Toffoli.

\textbf{MBU} (Jones 2013, applied to adder carry chains by Gidney 2018):
\begin{lstlisting}
1. Forward: ccx(a, b, out)       -> 4 T-gates (1 Toffoli)
2. Use 'out' in downstream logic
3. Uncompute: measure 'out' in X-basis (Hadamard, then Z-basis measure)
   Result is classical bit m in {0, 1}
   - If m = 0: done, free 'out'  -> 0 Toffoli
   - If m = 1: apply CZ(a, b)    -> 0 Toffoli (CZ is a Clifford gate)
Total uncompute cost: 0 Toffoli
\end{lstlisting}

\textbf{Why does this work?} When you compute \texttt{out = a AND b} into a clean \ket{0}
ancilla, the system is in state $|a, b, a\cdot b\rangle$. Measuring \texttt{out} in the
$X$-basis either succeeds cleanly ($m=0$) or applies a phase kick $(-1)$ to the
$|a=1, b=1\rangle$ component ($m=1$). The CZ gate corrects exactly that phase. After
correction, \texttt{out} is disentangled from $(a,b)$ and can be freed.

The key insight: \textbf{measurement collapses entanglement for free}. Coherent uncomputation
must undo entanglement gate-by-gate; measurement does it in one shot.

\textbf{The tradeoff}: The \texttt{out} qubit must remain live from when it's computed to
when it's measured (typically between the forward and reverse carry chains). This costs $+1$
peak qubit during that window.

\textbf{Practical effect on adders}: An $n$-bit ripple-carry adder normally costs $\approx 2n$
Toffoli ($n$ forward, $n$ reverse). With MBU on every carry-uncompute:
\[
  \text{Total adder cost} = n\ \text{Toffoli (forward only)} + 0\ \text{(measured uncompute)} = n\ \text{Toffoli}
\]
This halves the adder cost --- the most important constant-factor improvement in the SOTA.

\subsection{The Vent Dial}

\S\ref{sec:mbu} showed that MBU can \emph{uncompute} one carry bit for \textbf{0 Toffoli} (measure
it, then a free Clifford fix-up) instead of the 1 Toffoli a coherent uncompute would cost --- the
catch being that the carry qubit has to stay alive a little longer. Applying that choice to
\emph{one} carry in an adder is called a \textbf{vent}. The whole ``dial'' is just: \emph{how many of
an adder's carries do you choose to vent?}

\textbf{Where the $3n - 2$ comes from.} A reversible $n$-bit adder builds a chain of carry bits
climbing up the word (the \emph{forward} pass), uses the final carry, then unwinds that chain (the
\emph{reverse} pass) to return all the scratch qubits to \ket{0} (this compute--use--uncompute cycle
is the reversibility tax from \S\ref{sec:reversibility}). Tallying the Toffoli in that cycle for the
hybrid carry adder used here gives a fixed baseline of about $3n - 2$ Toffoli when \emph{nothing} is
vented. Each carry you vent deletes one Toffoli from the reverse (uncompute) half. So with $k$ vents:
\begin{align*}
  \text{Toffoli cost} &= 3n - 2 - k && \text{(one fewer Toffoli per vented carry)} \\
  \text{Peak qubits}  &= \text{baseline} + k && \text{(but each vented carry stays live a bit longer)}
\end{align*}
The dial is therefore perfectly linear and perfectly priced: \textbf{each vent $=$ $-1$ Toffoli,
$+1$ temporary peak qubit.}

\textbf{Why the qubit is only \emph{temporary} --- the ``forward$\leftrightarrow$reverse window.''} A
vented carry can't be freed the instant it's measured: its measurement outcome is still needed later,
when the reverse pass reaches that bit position to clean it up. So the qubit is occupied only from
when that carry is first computed (forward) until it is finally cleaned up (reverse) --- the stretch
in between is the ``forward$\leftrightarrow$reverse window.'' Outside that window the lane is free,
which is exactly why the $+k$ qubits don't add to the circuit's permanent footprint.

\textbf{When to vent --- it depends on \emph{where} the adder sits in the timeline.} Recall from
\S\ref{sec:measuring} that the circuit's qubit peak happens during a few specific phases (the
\emph{binder}), while other phases have \textbf{slack} --- spare qubits sitting idle. So:
\begin{itemize}
  \item \textbf{Off-peak adders} (running during a phase with slack): holding $+1$ qubit there does
    \emph{not} raise the \emph{global} peak, so venting is \textbf{pure profit} --- free Toffoli
    savings at zero score cost. Vent every carry you can.
  \item \textbf{On-peak adders} (running at the binder): every extra qubit pushes the peak --- and the
    score --- up. Vent only if the Toffoli saved is worth more than the qubit, judged by the exchange
    rate (\S\ref{sec:exchange}).
\end{itemize}

\textbf{How the SOTA automates it.} \texttt{trailmix\_ludicrous} doesn't choose by hand. It fixes a
circuit-wide qubit ceiling \texttt{TLM\_TARGET\_Q}, and for each adder call vents exactly
\texttt{TLM\_TARGET\_Q} $-$ \texttt{active\_qubits} carries --- i.e.\ it spends \emph{all} the spare
headroom beneath the ceiling on Toffoli savings, and not one qubit more. An adder already at the
ceiling gets $0$ vents (no room); an adder with lots of slack gets vented heavily. This is the
dynamic-headroom clamp (\S7.10) and the vent dial working together.

\begin{tcolorbox}[noteboxstyle, title=Gidney 2025 streaming vented adder (open)]
A newer construction (arXiv:2507.23079) that uses 2 clean $+$ $(n-2)$ \emph{dirty} (borrowed,
\S7.16) ancilla for $\approx 3n$ Toffoli. Partially implemented in \texttt{venting.rs} but
\textbf{currently leaks phase} --- the phase correction for the dirty ancilla is incomplete --- so it
is not yet in the SOTA. Expected $\approx n/4$ \tof savings per adder once fixed.
\end{tcolorbox}

\subsection{Conditional / Measured Replay (Average-Executed Only)}
\label{sec:replay}

\textbf{The idea}: If a comparator flag is \ket{0} on most inputs, measure the flag instead of
keeping it coherent. Then replay the dependent computation conditionally on the measurement outcome.

\begin{lstlisting}
1. Compute flag F into scratch ancilla (costs C Toffoli)
2. Measure F in X-basis -> classical bit m
3. If m = 0: dependent block is a no-op on this shot  -> 0 Toffoli from the block
4. If m = 1: replay the block classically conditioned on m  -> full cost
Average cost = C + Pr[F=1] * (block cost)
\end{lstlisting}

When $F$ fires on fraction $f$ of inputs: average-executed drops from $(\text{block cost})$ to
$f \times (\text{block cost})$. If $f = 0.05$, you pay only 5\% on average.

\textbf{The hard constraint}: This ONLY works on \textbf{self-uncomputing scratch} --- a comparator
you compute, read under the condition, and cleanly reverse. It does NOT work on data-controlled
operations or predicates reused downstream.

\subsection{Algebraic Fusion: Collapse Add/Negate Chains}
\label{sec:fusion}

Adjacent modular-arithmetic stages that are algebraically equivalent to a simpler single
operation can be collapsed to eliminate intermediate carry chains.

\textbf{Example} (\texttt{DIALOG\_FUSE\_C\_FORM}):
\begin{lstlisting}
Old:
  tx += 2*Qx    (mod add)
  tx = -tx      (mod negate)
  tx -= Qx      (mod sub)
  tx = -tx      (mod negate)

Algebraically (two negations cancel, constants combine):
  tx += 3*Qx    (one mod add)
\end{lstlisting}

\textbf{Example} (\texttt{DIALOG\_FUSE\_X\_RESTORE}):
\begin{lstlisting}
Old:
  tx = -tx      (mod negate)
  tx += Qx      (mod add)

Algebraically:
  tx = Qx - tx  (one operation)
\end{lstlisting}

Each removed negate/add/subtract eliminates a carry chain, a modular reduction fold, and a
cleanup phase --- saving significant Toffoli.

\subsection{Witness-First / Deferred-Output Dataflow}
\label{sec:witness}

Sometimes a cancellation only needs a \emph{witness relation}, not the final output value.
Computing the output early and then cleaning it is wasteful when the witness is cheaper.

\textbf{The 4352cfb Toffoli cut} ($-10.8$M Toffoli at fixed 980 qubits):

Key identity (from curve equation): $\texttt{new\_dy} = \lambda \cdot \texttt{new\_dx}$.
So to cancel $\lambda$, we do not need to compute \texttt{new\_y} first. Instead:
\begin{enumerate}
  \item Erase \texttt{dy} using: $dy = \lambda \cdot dx$ (subtract $\lambda \cdot dx$ from
    $dy \to$ zero)
  \item Reuse the zeroed \texttt{dy} register as scratch
  \item Compute \texttt{new\_x}
  \item Compute $\texttt{new\_dx} = ox - \texttt{new\_x}$
  \item Compute $\texttt{new\_dy} = \lambda \cdot \texttt{new\_dx}$ (the witness, not the output)
  \item Cancel $\lambda$ using $\texttt{new\_dy}/\texttt{new\_dx}$
  \item Materialize $\texttt{new\_y} = \texttt{new\_dy} - oy$ (only NOW compute the output)
\end{enumerate}

\textbf{Ablation}:
\begin{lstlisting}
zero-dy disabled, defer-y disabled:   110,227,695 emitted ops
defer-y materialization only:         105,573,887 emitted ops  (-4.6M)
both zero-dy + defer-y:                92,105,748 emitted ops  (-18.1M total)
\end{lstlisting}

\textbf{General principle}: Before materializing any output, ask: ``does the next operation NEED
the output, or just a function of it that's cheaper to compute directly?''

\subsection{Karatsuba Modular Square ($-22.4$ Million Toffoli)}

The largest single Toffoli reduction in the project's history.

\textbf{The schoolbook problem}: Computing $\lambda^2$ for a 256-bit number using schoolbook
multiplication costs $O(n^2)$ multiplications.

\textbf{Karatsuba}: For a 256-bit number split as $\lambda = \text{hi} \cdot 2^{128} + \text{lo}$:
\[
  \lambda^2 = \text{hi}^2 \cdot 2^{256}
            + \left[(\text{lo}+\text{hi})^2 - \text{lo}^2 - \text{hi}^2\right] \cdot 2^{128}
            + \text{lo}^2
\]
Instead of computing 4 products, compute only 3 squares: $\text{lo}^2$, $\text{hi}^2$,
$(\text{lo}+\text{hi})^2$. Reconstruct the cross-term via additions (cheap Clifford operations).

\textbf{Why $-22.4$M Toffoli}: The squaring operation runs at EVERY GCD step (258 steps $\times$
2 passes). Any $O(n^2) \to O(n^{1.585})$ improvement compounds massively across all those calls.

\textbf{Density-neutral}: The Karatsuba identity is a pure algebraic identity, correct for all
integers.

\textbf{Recursive Karatsuba}: Tested on 128-bit sub-squares $\to$ net Toffoli LOSS (the overhead
of the additions outweighs the sub-square savings at that level). Dead end.

\textbf{The space/Toffoli tension (it cuts both ways).} Karatsuba saves Toffoli but \emph{needs
scratch} for the three sub-products and the $(lo+hi)^2$ term. When the square sits at the qubit peak,
the \textbf{opposite} move can win: drop back to a plain \textbf{schoolbook} square, which needs no
middle-term scratch (fewer qubits) at the cost of more Toffoli. The SOTA toggles this per context
(\texttt{ROUND84} ``schoolbook x-tail'') --- Karatsuba where Toffoli is binding, schoolbook where
\emph{qubits} are. Same exchange-rate decision as everywhere else (\S\ref{sec:exchange}): which axis
is binding decides which square you want.

\subsection{NAF Recoding of Constants}

\textbf{Non-Adjacent Form (NAF)} uses signed digits $\{-1, 0, +1\}$ to represent constants with
fewer non-zero digits:
\[
  977 = 2^{10} - 2^{5} - 2^{4} + 1 \quad \text{(4 signed terms vs 6 binary terms)}
\]

In NAF, no two adjacent digits are both non-zero, and the number of non-zero digits is minimal.
For $977$: 4 terms instead of 6, saving $\approx33\%$ of the Solinas fold cost.

\textbf{Density-neutral}: The identity $977 = 1024 - 32 - 16 + 1$ is a mathematical fact, correct
for all inputs.

\subsection{Dead-CCX Elimination (Empirical Drop vs Structural Skip)}
\label{sec:deadccx}

Some CCX gates contribute nothing to the answer --- either their control is always zero, or their net
effect cancels --- yet the scorer still \emph{charges} them. Deleting such ``dead'' gates lowers the
average-executed Toffoli. There are two ways to find them, with \textbf{very different correctness
profiles}, and the difference is the clearest example of the overfitting trap in the whole project.

\textbf{(1) Empirical dead-CCX drop (the island-overfit lever).} Run the circuit over a large sample
of inputs ($\sim$9 million reachable EC points, or many Fiat--Shamir test sets); record which CCX
gates \textbf{never flip their target} --- i.e.\ are \emph{inert-but-charged} on the sample. Collect
their positions into a \texttt{.idx} file (e.g.\ 13{,}873 op indices) and filter them out:

\begin{lstlisting}
ops = ops.filter(|(i, _)| !drop_set.contains(i))   // delete the inert-but-charged CCX
\end{lstlisting}

This was the engine behind the SOTA average-Toffoli for a long time (a single drop was worth
$\sim$13{,}000$+$ avg-Toffoli). \textbf{But it is distribution-exact / island-overfit, and dangerous
for three reasons:}
\begin{itemize}
  \item \textbf{It deletes gates that \emph{could} fire on unsampled inputs.} The result is ``clean
    on its island'' ($0/0/0$ over the hunted 9024) but \textbf{not} provably correct off-island ---
    a gate inert on every sampled input may be live on one you never tried, and the deletion then
    silently corrupts the output.
  \item \textbf{The indices are absolute positions in one exact op stream.} \emph{Any} structural
    change --- a clamp, a knob, a fold toggle, even reordering --- shifts every later gate onto the
    wrong index, so a stale list deletes \emph{live} gates ($9024/141/141$ all-shots breakage). Every
    variant must ship its own freshly-screened list and freshly-hunted nonce.
  \item \textbf{It does not improve the construction.} A genuinely better circuit would not
    \emph{emit} these gates; the drop just hides them from the scorer, and it is the first thing
    invalidated by any real structural improvement. Treat it as a sophisticated nonce grind --- worth
    knowing it exists (it is why the SOTA avg-T looked low), low-value to reimplement except to
    squeeze a \emph{frozen} final artifact.
\end{itemize}

\textbf{Iterated (two-pass) empirical drop} (\texttt{DROP\_DEAD\_ROBUST\_SECOND}): the drop is
\textbf{iterable}. After the first drop reshapes the stream, re-screen the \emph{already-dropped}
stream --- a second screen finds gates newly (or only now detectably) dead --- drop a second, smaller
list, and re-hunt the nonce with \emph{both} drops on. Subset-monotone with diminishing returns (2nd
pass $\sim$$-2{,}411$ avgT at 1153q). All overfitting caveats above apply at every pass.

\textbf{(2) Structural dead-CCX skip (the sound, density-neutral lever --- CURRENT SOTA).} Instead of
\emph{sampling}, \textbf{prove} a CCX never fires on \emph{any} input from circuit structure alone
(known-constant carries, exact-remainder conditions, call-site bit patterns). Named predicates in
\texttt{trailmix\_ludicrous}: \texttt{TLM\_FFG\_SKIP\_STRUCTURAL\_DEAD\_*},
\texttt{TLM\_CUCCARO\_SKIP\_STRUCTURAL\_DEAD\_*}, \texttt{TLM\_COMPARE\_SKIP\_STRUCTURAL\_DEAD\_*}.
Keyed by call-site and bit-index, not sampled data, so they are \textbf{density-neutral} (correct on
all inputs, stackable freely, no \texttt{.idx}, no re-screen).

\textbf{The 6dafa07 pivot.} The SOTA \emph{replaced} the empirical drop entirely with structural
skipping --- which turned out to be \textbf{both sound and lower-Toffoli}. The general lesson in
miniature: when an island-exact lever is doing real work, look for the \emph{structural} reason those
gates are dead and skip them provably. If the gates truly never fire, you can usually prove it --- and
the overfitting risk then vanishes for free. (See \S\ref{sec:density} for the value-exact vs
island-exact framing.)

\subsection{Classical Constant Folding}

When one operand of an arithmetic operation is a \textbf{classical constant} (known at compile
time), the operation can often be implemented with 0 Toffoli gates.

\textbf{The mechanism}: A \texttt{BitId} register holds a constant pattern as compile-time
information. Operations with one classical operand:
\begin{itemize}
  \item \texttt{qubit XOR 0} = identity (no gate)
  \item \texttt{qubit XOR 1} = Pauli $X$ (free Clifford gate)
  \item \texttt{qubit + classical\_constant mod p} = a series of CNOT-like operations --- 0 Toffoli.
\end{itemize}

\textbf{In trailmix\_ludicrous}: $Q = (Q_x, Q_y)$ is classical, so \texttt{coord\_add3x(circ, x2,
ox, qubits)} (\textit{x2 $\mathrel{+}$= 3$\cdot Q_x$ mod p}) costs \textbf{0 Toffoli}.

\subsection{Converged-Tail Gate Elision}

In the last $K$ iterations of the 530-step GCD, the algorithm has converged and the apply-phase
\texttt{cswap} has a control that is deterministically 0 on all but rare inputs. Skip the cswap
for those last $K$ iterations.

\textbf{Cost}: Island-exact. Requires re-hunting the tail nonce after applying.

\subsection{GAP\_J2 Comparator Window Narrowing ($-1.33$M Toffoli)}

The swap decision comparator at each GCD step checks whether $|u| > |v|$. A full 256-bit
comparison would be expensive, but the GCD ensures that $u$ and $v$ typically differ somewhere
in their top bits.

\textbf{The schedule \texttt{GAP\_J2[i]}} (258-element table) sets the comparator window width
per GCD step. The comparator scans only the top $\texttt{GAP\_J2[i]}$ bits. A 22-line edit to the
table to shave $\approx1$ bit per step across 258 steps:
\[
  \text{Toffoli saved} = 1\ \text{bit} \times 516\ \text{comparator calls} \approx 1.33\text{M Toffoli}
\]
\textbf{Island-exact} --- requires nonce hunting.

\subsection{Ancilla-Free Majority}
\label{sec:majority}

Standard carry logic uses a 3-input majority gate with 3 Toffoli and an ancilla. The identity:
\[
  \text{maj}(a, k, c) = c \oplus ((a \oplus c) \wedge (k \oplus c))
\]
implements the same majority with \textbf{2 Toffoli and no ancilla} (XOR $=$ CNOT, AND $=$ CCX).
Peak-neutral, value-identical, pure count reduction wherever majority gates appear.

\subsection{Exact-Adder Recovery}

Counter-intuitive: Sometimes a truncated adder costs \textbf{more} Toffoli than the exact
full-carry adder, because the truncation requires an expensive correction to handle the rare
carry-escape cases.

\textbf{The rule}: If the correction overhead of a truncated adder exceeds its savings from
dropping bits, revert to the exact adder. Named: \texttt{DIALOG\_GCD\_APPLY\_FINAL\_TOPCLEAN=0}
$\to$ $\approx{-2{,}597}$ avg-Toffoli, value-exact.

\subsection{Skip Redundant Final Cleanup}
\label{sec:skipclean}

A folded/structured computation sometimes makes a final ``top-clean'' pass redundant --- the
structure ensures the cleanup bits are already zero by construction. Dropping the cleanup removes
its Toffoli.

\begin{tcolorbox}[noteboxstyle, title=Highest-risk move]
If the cleanup is not actually redundant, you get silent phase/classical dirt on some inputs.
Always pair with strong \cls/\pha/\anc{} validation before enabling.
\end{tcolorbox}

\subsection{Off-Peak Loosening (The Dual Move)}

When a phase falls \textbf{off} the qubit peak (another phase is now the binder), you can WIDEN
that phase's arithmetic to recover Toffoli for free:
\begin{lstlisting}
Phase A: was binder at 1170q -> locally reduced to 1162q
Phase B: now binder at 1170q (unchanged)
Action: widen Phase A's carry width back toward the old value
Effect: Phase A's Toffoli drops; Peak stays at 1170 (Phase B is binding)
\end{lstlisting}

\textbf{The rule}: Shrinking an off-peak phase is a pure-Toffoli play. Widening an off-peak phase
is a free Toffoli refund. Govern both on gates alone, not qubits.

\subsection{Constprop: Automated Static Gate Elimination}

The constraint-propagation pass (\texttt{CONSTPROP\_MAX\_ITERS}) analyzes the circuit statically
and eliminates:
\begin{itemize}
  \item \textbf{Provably-constant CCX gates}: control always 0 $\to$ remove; control always 1
    $\to$ replace with CNOT
  \item \textbf{Inverse-pair cancellation}: consecutive \texttt{CCX(a,b,c)} gates that cancel
    algebraically
  \item \textbf{Affine simplification}: sequences simplifying to an affine function, implementable
    with fewer gates
\end{itemize}

Setting \texttt{CONSTPROP\_MAX\_ITERS=256} (vs default 16) runs to a deeper fixpoint, finding
more opportunities: $-377$k Toffoli in one version.

\subsection{Quadratic-Cascade $\to$ Controlled-Increment (cinc)}
\label{sec:cinc}

\textbf{Replace an $O(t^2)$ gadget with an $O(t)$ one when the math allows.} A recurring pattern is
a ``carry-into-tail'' or prefix-propagation built as a \textbf{quadratic cascade of
multi-controlled-X gates} (\lstinline{mcx_clean_k} over a window $[nv, L)$ --- every position
controls on all positions below it, giving $O(t^2)$ MCX work).

When the cascade is computing a \emph{carry propagation} (the standard ``increment if all low bits
are 1'' pattern), it is exactly a \textbf{controlled increment}, and there is a cheaper exact
primitive for that: the \textbf{Khattar--Gidney controlled-increment} \texttt{cinc}
(arXiv:2407.17966), which uses only $\log^* n$ clean ancilla and $O(n)$ Toffoli instead of the
quadratic cascade.

\begin{lstlisting}
Old: clean-tail fold carry = mcx_clean_k prefix cascade over [nv, L)   -> O(t^2) MCX
New: clean-tail fold carry = cinc_khattar_gidney (one controlled +1)   -> O(t), log*-ancilla
\end{lstlisting}

\textbf{Value-exact and density-neutral}: an exact functional replacement (same truth table), so it
does not touch the island distribution --- it can be stacked freely without re-hunting a nonce.
Named \texttt{TLM\_FOLD\_TAIL\_CINC}; worth $\sim$$-5{,}175$ avgT alone. This is the \textbf{``safe''
Toffoli lever class}: unlike dead-CCX deletion (\S\ref{sec:deadccx}, distribution-exact), a cinc swap
is correct on all inputs --- prefer this class when choosing what to stack first.

The general principle: \textbf{whenever a gadget's cost is quadratic in a window width, ask whether
it is secretly a standard primitive (increment, compare, prefix-AND) with a known linear
construction.}

\subsection{Pseudo-Mersenne (Solinas) Modular Arithmetic}
\label{sec:pmarith}

This is the workhorse that makes \emph{every} modular operation cheap. The math background is in
\S\ref{sec:pmmath}: because $p = 2^{256} - c$ with $c = 2^{32} + 977$ only 33 bits,
$2^{256} \equiv c \pmod p$, so a reduction is a cheap \textbf{fold} of the top bits into the bottom,
not a division. Here is how that identity becomes actual reversible gadgets (Schrottenloher
Algs 7/10/11; trailmix \lstinline{pm_prims.rs} / \lstinline{arith.rs}). The constant $c$ is called
$f$ in the code.

\textbf{Modular doubling --- $2x \bmod p$ (Alg 7).} Shift $x$ left by one bit (now up to 257 bits).
The part that spilled past bit 256 has weight $2^{256} \equiv c$, so it folds back as $+c$:
\begin{lstlisting}
2x mod p:
  1. shift-left by 1                       (a relabel -- 0 Toffoli)
  2. if the top bit overflowed: add c into the LOW limbs   (controlled +f, short)
  3. one CX to clear the overflow ancilla
\end{lstlisting}

\textbf{Modular addition --- $x + y \bmod p$ (Alg 10).} $x + y$ overflows $2^{256}$ by at most one
bit; on overflow, fold it as $+c$ in the low limbs, then erase the overflow flag with a compare:
\begin{lstlisting}
x + y mod p:
  1. add y into x                          (a Cuccaro/Gidney adder)
  2. if overflow (carry-out = 1): add c into the low limbs
  3. erase the overflow ancilla via a y<x comparison
\end{lstlisting}
Alg 11 additionally handles the degenerate $x + y = p$ case (all high bits equal) --- needed only in
the GCD's first $\approx c\sqrt{n}$ ``empty'' padding iterations, so the expensive variant is routed
\emph{only} there and the cheap Alg-10 used everywhere else.

\textbf{The two approximations that make it cheap (and their dials).} Doing these exactly would still
need full-width compares and adds. Two structural facts let you truncate --- because $c$ is small:
\begin{itemize}
  \item \textbf{$+f$ touches only the low limbs (the \texttt{lsbs} dial).} Folding $+c$ perturbs only
    the bottom $\sim$33 bits plus a short carry ripple, so the $+f$ adder runs over a window of width
    \texttt{lsbs} (trailmix uses $\sim$54--63). Too narrow misses a long ripple; per-call failure
    $\approx 2^{-\text{padding}} \approx 2^{-30}$ at their settings. \emph{Shipped as}
    \texttt{*\_CARRY\_TRUNC\_W} / \texttt{ROUND84\_INPLACE\_SOLINAS\_FOLD} / the low-54-bit $+f$ fold
    (\S\ref{sec:trailmix}).
  \item \textbf{Overflow/degeneracy decided from only the top bits (the \texttt{msbs} dial).} Whether
    $x+y \ge p$ (or $2x \ge p$) is dominated by the high bits, so the compare scans only the top
    \texttt{msbs} bits (the paper uses 40--50). \emph{Shipped as} \texttt{DIALOG\_GCD\_COMPARE\_BITS}
    / \texttt{*\_APPLY\_CLEAN\_COMPARE\_BITS}.
\end{itemize}
Both are \textbf{graded / island-exact} levers (\S\ref{sec:qubitloop}, \S\ref{sec:density}):
tightening a notch is value-exact only on a searched nonce-island, because it makes a $2^{-k}$-rare
fraction of inputs reduce wrong --- exactly the binary-vs-graded line of \S\ref{sec:qubitloop}, tuned
against the 9024-shot $0/0/0$ target.

\textbf{Borrowed-dirty $+f$ window (\lstinline{gidney_cadd_f_window}).} The $+f$ add needs carry
scratch; instead of fresh ancilla it \textbf{borrows the register's own idle high bits} ($\sim$3
clean ancilla instead of $\sim$10), restoring them after --- gate-hosting (\S\ref{sec:gatehost})
applied to the fold. The reduction adds $\sim$0 clean qubits at the peak, value/phase-identical to
the vented path.

\textbf{Why it matters.} A generic-prime reduction is a full multiprecision divide
(Barrett/Montgomery, $\sim$n-bit work) on \emph{every} add and double, and the GCD/B\'ezout
reconstruction does thousands of them. Pseudo-Mersenne turns each into a short $+f$ fold $+$ a
top-bits compare --- the entire secp256k1 advantage (constant-adder 15\% vs 34\%, modular-double 8\%
vs 24\%, \S\ref{sec:pmmath}). \textbf{No Montgomery form is used} --- both the paper and trailmix
keep plain integer representation and fold directly. NAF recoding (\S9.7) is the next layer: it makes
the fold constant $c$ itself cheaper (4 signed terms instead of 6).

\subsection{Merging Adjacent Controlled-Swaps}
\label{sec:cswapmerge}

The GCD loop swaps its two operand registers conditionally at almost every step (the ``is
$|u| > |v|$?'' decision swaps $u$ and $v$, and likewise their B\'ezout cofactors). A controlled swap
of two $n$-bit registers costs $n$ Toffoli (one \texttt{cswap} per bit). Across $\sim$530 steps
$\times$ two register pairs, these swaps are a large slice of the Toffoli budget.

\textbf{The identity that merges them.} Two controlled swaps of the \emph{same pair of registers},
controlled by bits $p$ and $q$, compose into a \emph{single} controlled swap controlled by
$p \oplus q$:
\[
  \texttt{cswap}(p, A, B)\cdot\texttt{cswap}(q, A, B) \;=\; \texttt{cswap}(p \oplus q, A, B).
\]
(Swapping under $p$, then again under $q$, swaps exactly when an odd number of the two controls fired
--- i.e.\ when $p \oplus q = 1$.) So wherever two swap decisions on the same registers sit next to
each other --- e.g.\ the cofactor swap at the \emph{end} of one GCD step and the operand swap at the
\emph{start} of the next --- you replace \textbf{two} $n$-Toffoli swaps with \textbf{one}, after
spending a single cheap CNOT to compute $p \oplus q$ into a ``frame-parity'' qubit that carries the
accumulated swap state across the boundary.

This is \textbf{algebraic fusion (\S\ref{sec:fusion}) specialized to swaps}, and it is value-exact (a
pure identity, all inputs). It was the single biggest Toffoli win of the early era --- the Kaliski
\texttt{step9}$\circ$\texttt{step3} boundary merge, \textbf{$-274{,}000$ Toffoli, peak-neutral}. The
general lesson: \emph{controlled swaps compose through XOR of their controls} --- look for adjacent
swaps on the same lanes and collapse them.

% ===========================================================================
\section{The Qubit $\leftrightarrow$ Toffoli Exchange Rate and the Product Race}
\label{sec:exchange}
% ===========================================================================

\subsection{The fundamental arithmetic}

At the current SOTA ($1152\text{q} \times 1{,}364{,}230\ \text{avgT} = 1{,}571{,}592{,}960\ \text{score}$):
\[
  \text{break-even rate} = \frac{\text{avgT}}{\text{peak\_q}} = \frac{1{,}364{,}230}{1152} \approx 1{,}184\ \text{Toffoli per qubit}
\]
A lever that saves 1 qubit at the cost of fewer than 1,184 Toffoli \textbf{improves the score}.
A lever that costs more than 1,184 Toffoli per qubit saved \textbf{worsens the score} --- even
though it reduced qubits.

This rate \textbf{changes with every optimization}:
\begin{itemize}
  \item After a large Toffoli reduction: the rate drops (qubits become more valuable). Previously-marginal qubit cuts are now worth trying.
  \item After a large qubit reduction: the rate rises. Previously-marginal Toffoli cuts might now dominate.
\end{itemize}

\subsection{The product-race caveat}

A qubit drop that clears break-even can STILL lose the product race if the Toffoli-grinding
track moves faster in parallel.

\textbf{Concrete example} (June 2026):
A 1157q clamp cost $\approx1{,}127$ Toffoli/qubit --- below break-even at landing. But the 1159q
Toffoli-grind reached $1{,}378{,}242$ avgT. Score comparison:
\[
  1159 \times 1{,}378{,}242 = 1{,}597{,}506{,}378 < 1157 \times 1{,}380{,}890 = 1{,}598{,}290{,}330
\]
The 2-qubit savings was overtaken by the parallel Toffoli-cutting track.

\textbf{Lesson}: Always run both tracks and compare products. Never minimize qubit count in isolation.

\subsection{The shelved-lever rule}

A qubit lever that \textbf{failed break-even on one Toffoli base} can become favorable on a
cheaper Toffoli base.

\textbf{Example}: The dynamic headroom clamp (\texttt{TLM\_TARGET\_Q}) cost $\approx1{,}514$
Toffoli/qubit on the schoolbook-square base. After Karatsuba reduced the square cost, the same
clamp cost only $\approx20$ Toffoli/qubit on the new base --- well below break-even. It won.

\textbf{Rule}: After any major arithmetic win (Karatsuba, algebraic fusion, structural rewrite),
re-test \textbf{all previously-shelved qubit levers} at the new exchange rate.

\subsection{Different peak layers have different exchange rates}

At any peak, there are typically two layers:

\begin{enumerate}
  \item \textbf{Persistent floor}: long-lived state held across the whole peak (transcript,
    passenger values, GCD state). Often \textbf{cheap} to reduce --- a denser codec or transcript
    streaming costs few gates to save many qubits.
  \item \textbf{Transient working registers}: the current step's carry bits, arithmetic scratch.
    Often \textbf{expensive} to reduce --- freeing them requires recompute/dirty-borrow, adding
    many Toffoli.
\end{enumerate}

Counter-intuitively: \textbf{compress the transcript before freeing the scratch}. The transcript
looks ``irreducible'' (structured data) but is often the cheap lever. The scratch looks
``wasteful'' but is often the expensive lever.

% ===========================================================================
\section{Density-Neutral vs Island-Exact --- Correctness Regimes}
\label{sec:density}
% ===========================================================================

\subsection{The fundamental distinction}

\begin{itemize}
  \item \textbf{Density-neutral (value-exact)}: The circuit produces the same output on ALL
    inputs. The optimization is grounded in mathematical identity or circuit structure.
  \item \textbf{Island-exact (distribution-specific)}: The circuit is correct only on the
    specific 9024 test inputs. On other inputs, it may give wrong answers.
\end{itemize}

When you apply an island-exact optimization, you need to find a new \textbf{nonce} --- a 48-byte
seed for the test input generator --- such that all 9024 generated inputs happen to be in the
safe zone. Stacking multiple island-exact optimizations makes the safe zone smaller.

\subsection{A common island-exact lever: active-iteration trimming}

The GCD inversion is run for a fixed number of steps chosen to cover the \emph{worst-case} input. But
the worst case is rare, so you can run \textbf{fewer} iterations
(\texttt{DIALOG\_GCD\_ACTIVE\_ITERATIONS}, e.g.\ $399\to393$) and accept that a $\sim2^{-k}$-rare
fraction of inputs won't have fully converged. This is a pure island-exact lever --- it drops whole
divsteps (qubits \emph{and} Toffoli) but makes the circuit wrong on those rare inputs, so it must be
paired with a nonce hunt. It sits alongside carry-truncation (\S7.4) and comparator-narrowing
(\S9.11/\S9.18) as a ``graded'' knob (\S\ref{sec:qubitloop}).

\begin{tcolorbox}[noteboxstyle, title={The line between island-exact and cheating}]
Island-exact levers are legitimate \emph{because the 9024 validated inputs are the actual scored
distribution} and the failures they introduce are genuinely rare and structurally understood. There
is a real line, though: deliberately truncating the circuit so it is wrong on a \emph{large} fraction
of inputs and then re-rolling the nonce purely to \emph{hide} those failures is verifier-gaming, not
optimization --- the early history contains at least one flagged ``red-team'' example (a 396-iteration
build wrong on $\sim$$3\times10^{-4}$ of inputs, shipped with a comment admitting it). The community
norm: an island-exact cut should be a structurally-sound approximation whose error you can
\emph{bound and explain}, not a way to launder a broken circuit past the gate.
\end{tcolorbox}

\subsection{The 6dafa07 pivot: empirical $\to$ structural dead-CCX}

Before the current SOTA, the circuit used \textbf{empirical dead-CCX elimination}: sample 9M
inputs, find CCX gates that never fire, skip them. Island-exact.

The \texttt{6dafa07} commit replaced this with \textbf{structural dead-CCX skipping}: prove from
circuit structure which CCX gates can never fire on \emph{any} input. These predicates are keyed
by call-site and bit-index, not sampled data. Density-neutral.

\subsection{The correctness check triple: \cls/\pha/\anc}

Every validation run checks three error types:
\begin{itemize}
  \item \cls: classical mismatch (the output value is wrong on some input)
  \item \pha: phase error (the output is correct but entangled with extra phase --- kills interference)
  \item \anc: ancilla mismatch (a qubit failed to uncompute back to \ket{0})
\end{itemize}

A correct circuit shows \texttt{0/0/0}. Measurement-based uncompute fails in \pha{} (which value
checks miss). Truncation failures typically show in \cls. Missed uncompute shows in \anc.

% ===========================================================================
\section{Pebbling Theory --- The Formal Foundation}
\label{sec:pebbling}
% ===========================================================================

\subsection{The reversible pebble game}

The \textbf{pebble game} models space-time tradeoffs in reversible computation. The computation is
a directed graph (nodes $=$ values, edges $=$ dependencies). A ``pebble'' $=$ a qubit holding that
value. Peak qubits $=$ maximum pebbles simultaneously.

\subsection{Bennett pebbling: classical reversibility is expensive}

Bennett (1989) showed that reversibly simulating a $T$-step computation using $S$ extra pebbles
costs:
\[
  \text{Time} = \varepsilon \cdot 2^{1/\varepsilon} \cdot T^{1+\varepsilon} / S^\varepsilon \quad \text{(for any $\varepsilon > 0$)}
\]
The hidden constant $c(\varepsilon) = \varepsilon \cdot 2^{1/\varepsilon}$ grows \textbf{exponentially}
as $\varepsilon \to 0$.

\textbf{Concrete example}: Halving the GCD transcript ($S$ halved, $\varepsilon = 1$):
\[
  \text{Time blowup} \approx 1 \times 2 \times 530^2 / 370 \approx 1{,}518\ \text{extra GCD passes}
\]
At $\approx2{,}500$ Toffoli per pass: 3.8M extra Toffoli for $\approx370$q savings.
Break-even needs: $3.8\text{M} / 1{,}184 \approx 3{,}200$ qubits. Actual saving: 370. Not worth it.

\subsection{Quantum (spooky) pebbling: exponentially better}

\textbf{Gidney (2019)} introduced \textbf{spooky pebbling}: also allow removing a pebble by
measuring it in the $X$-basis, recording the measurement outcome as a classical bit. When the
value is needed again, recompute it and apply a classically-conditioned phase correction.

\textbf{Kornerup, Sadun, Soloveichik (2021)} proved tight bounds:
\begin{align*}
  \text{Quantum:}   &\quad T_{\text{quantum}} = O(T/\varepsilon)           \quad [\text{polynomial constant}] \\
  \text{Classical:} &\quad T_{\text{classical}} = O(2^{1/\varepsilon} \cdot T) \quad [\text{exponential constant}]
\end{align*}
The quantum advantage is exponential. This is the theoretical backbone of why MBU and
measurement-based techniques dominate.

\subsection{Parallel spooky pebbling (Kahanamoku-Meyer et al.\ 2025)}

For a length-$\ell$ sequential computation chain:
\[
  \text{qubits needed} = 2.47 \times \log(\ell) \quad \text{at depth } 2\ell
\]

For our 530-step GCD: $2.47 \times \log(530) \approx 23$ qubits. Our current GCD state machine
uses $\approx22$ qubits --- already near the theoretical minimum.

\subsection{How the optimization techniques \emph{are} pebble-game moves}

This theory is not background decoration --- \textbf{every qubit-freeing technique in
\S\ref{sec:qubittech}--\S\ref{sec:toftech} is literally a move in one of these two pebble games.}
Naming the move tells you its exact cost and risk:

\begin{center}
\small
\begin{tabular}{p{3.4cm}p{5.2cm}p{2.4cm}p{2.6cm}}
\toprule
\textbf{Technique} & \textbf{Pebble-game move} & \textbf{Which game} & \textbf{Cost to un-pebble} \\
\midrule
Live-range hole (\S7.1) & remove a pebble by \textbf{uncompute}, re-place later by \textbf{recompute} & Bennett & full Toffoli, \emph{twice} \\
Free-and-recompute cy0 (\S7.5) & same, on one cheap pebble & Bennett & $\approx$0 (a few CNOTs) \\
MBU / carry vent (\S9.1--B) & remove a carry pebble by \textbf{X-measurement $+$ CZ} & spooky & \textbf{0 Toffoli} \\
Conditional replay (\S9.3) & measure the pebble; recompute only on the firing branch & spooky & $\Pr[\text{fire}] \times$ cost \\
HMR ghosting (\S7.14) & measure, leave a \textbf{ghost}, recompute $+$ phase-fix & spooky & one recompute \\
Transcript compression (\S7.8) & \textbf{don't pebble it} --- shrink the data so fewer pebbles persist & (out of the game) & codec only \\
Parallel GCD schedule (\S\ref{sec:trailmix}) & optimal pebble \emph{placement} over the 530-step chain & parallel spooky & --- (already optimal) \\
\bottomrule
\end{tabular}
\end{center}

\textbf{The Bennett$\to$spooky upgrade is the single biggest constant-factor lever.} A live-range
hole (\S7.1) is a \emph{Bennett} move: to free a pebble you uncompute it (paying its Toffoli) and
later recompute it (paying again) --- the classical-reversible game, where un-pebbling is expensive.
MBU (\S9.1) is the \emph{spooky} move on the same kind of pebble: you remove it by \textbf{measuring}
it, and a measurement costs no Toffoli. That is exactly why an $n$-bit adder drops from $\sim$2n
Toffoli (compute $+$ coherent uncompute) to $\sim$n (compute $+$ measured uncompute) --- \textbf{MBU
is the pebble game's quantum advantage cashed out on every carry bit.}

\textbf{Pebbling theory also tells you \emph{which} qubits to attack.} The recompute-to-free move
only pays when $\texttt{recompute\_Toffoli} < \texttt{qubits\_freed} \times \texttt{break\_even}$
(\S\ref{sec:exchange}). The two extremes:
\begin{itemize}
  \item \textbf{The transcript is firmly in the ``never recompute'' regime.} Bennett's
    exponentially-growing constant is why: re-deriving it costs $\sim$3.8M Toffoli to free $\sim$370
    qubits ($\sim$8$\times$ over break-even). So take it \emph{out} of the game by \textbf{compressing
    the data} (\S7.8), not by recomputing it.
  \item \textbf{A scratch pebble like \texttt{cy0} is deep in the ``always recompute'' regime.} It is
    a cheap function of live qubits, so recompute costs $\sim$0 Toffoli --- free it without hesitation
    (\S7.5).
\end{itemize}

And the parallel-spooky bound says the \textbf{GCD control chain is already done} ($\approx$22q vs
the $2.47\cdot\log(530) \approx 23$q floor). So the theory aims every remaining qubit lever at the
\textbf{data pebbles} --- transcript, passengers, coordinate registers --- exactly where \S7 and
\S\ref{sec:trailmix} spend their effort. In short: \emph{recompute the cheap pebbles (Bennett), measure
the carry pebbles (spooky), compress the expensive data instead of pebbling it, and stop optimizing
the control chain.}

% ===========================================================================
\section{The SOTA Circuit: How trailmix\_ludicrous Works}
\label{sec:trailmix}
% ===========================================================================

\subsection{The decisive architectural choice: Q is classical}

The second point $Q = (Q_x, Q_y)$ is kept as classical \texttt{BitId} values, not quantum
registers. Materialized into a transient quantum temp \emph{only} at off-peak coordinate steps,
then freed. This eliminates 512 qubits from the peak.

\subsection{Two GCD passes sharing one tape}

The affine point addition needs:
\begin{enumerate}
  \item $\lambda = dy \cdot dx^{-1} \bmod p$ --- inverse GCD (drive $(p, dx) \to (1, 0)$)
  \item $\text{new\_y\_component} = \lambda \cdot (P_x - P_x') \bmod p$ --- forward GCD
\end{enumerate}

Both share \textbf{one} modular-inverse engine and \textbf{one} dialog tape (transcript). This
avoids having two separate inversion circuits, saving both qubits (one shared tape) and Toffoli
(shared codec overhead).

\subsection{The base-5 codec: 772 $\to$ 603 qubits live}

At each of 257 divstep steps, the algorithm records $(\texttt{subtracted}, \texttt{swap}, s_2)$ ---
3 bits raw. But only 5 of 8 possible 3-bit patterns are reachable. Triple windows have $5^3 = 125$
reachable states, encoded in 7 bits instead of 9:
\begin{lstlisting}
Layout: 1 Step0 (2b) + 85 Triples (7b each) + 1 Pair tail (5b) + 1 t1 flag (1b) = 603 qubits
\end{lstlisting}

\subsection{The vented-adder zoo}

Every adder in \texttt{trailmix\_ludicrous} vents its carry-uncompute via MBU:
\begin{lstlisting}[language=Rust]
hmr(carry, bit);       // X-basis measurement, 0 Toffoli
cz_if_bit(a, b, bit);  // phase correction conditioned on measurement
\end{lstlisting}

The adder family dispatches per-call by baked schedule tables. Each call's vent budget: exactly
\texttt{TLM\_TARGET\_Q} $-$ \texttt{active\_qubits}, spending every spare qubit below the ceiling
on Toffoli savings.

\subsection{The deliberate island-exact components}

trailmix\_ludicrous includes calibrated island-exact approximations:
\begin{itemize}
  \item \textbf{Low-54-bit +f fold} (\texttt{LSBS = 54}): reduction folds touch only low 54 bits;
    carry beyond bit 53 is dropped with probability $\approx 2^{-21}$ per fire.
  \item \textbf{Narrow-top-window comparators} (\texttt{MSBS = PAD = 21}): swap predicates scan
    only the top 21 bits. Mis-decision probability $\approx 2^{-21}$ per call.
\end{itemize}

\subsection{Where 1152q comes from}

At the SOTA (\texttt{d44cad3}):
\begin{itemize}
  \item \textbf{603q}: compressed base-5 GCD transcript (live during both GCD passes)
  \item \textbf{512q}: two 256-bit coordinate registers (quantum $P_x$, $P_y$)
  \item \textbf{$\approx37$q}: transient carry/vent ancilla at the peak-binding fold operations
\end{itemize}

% ===========================================================================
\section{The Three Tracks: Score, Low-Qubit, and Pareto Frontier}
\label{sec:history}
% ===========================================================================

The challenge has three distinct objectives (\S\ref{sec:twocomps}), and each has its own history and
its own winning constructions. This section walks all three: the \textbf{score track} (minimize
Q$\times$T), the \textbf{low-qubit track} (minimize Q alone), and the \textbf{Pareto-frontier track}
(map the clean (Q, T) curve).

\subsection{The score track (Q$\times$T): the 2715q$\to$1152q journey}

The score track has three clearly distinct eras, each defined by its \textbf{inversion engine}: the
Kaliski pre-history, the dialog-GCD frontier (which stalled at a 1168q wall for days), and the
\emph{ludicrous} breakthrough that broke it and drove to today's SOTA.

\subsubsection{Era 1 --- Pre-history: 2715q $\to$ 1698q (the Roetteler/Kaliski era)}

The challenge \emph{started} at a baseline of \textbf{2{,}715 qubits $\times$ $\sim$3.96M Toffoli}: a
\textbf{Roetteler-style two-Kaliski affine} point addition --- affine coordinates (not projective),
the modular inverse done by \emph{two} runs of the Kaliski binary-GCD algorithm, direct/Solinas
reduction (no Montgomery), with the prime and $Q$ known as compile-time constants. The inverse was
$\sim$80--90\% of the Toffoli (the bottleneck \S\ref{sec:circuit} names), treated as algorithmically
fixed --- so the early drops were \emph{all} live-range and hosting tricks, the same methods in this
primer in cruder form:

\begin{itemize}
  \item \textbf{2708 $\to$ $\sim$2002 (scratch-packing on a fixed inverse):} the
    $\texttt{cswap}(p\oplus q)$ swap-merge (\S\ref{sec:cswapmerge}, $-274$k Toffoli), borrowing
    \textbf{provably-\ket{0} high bits} of the shrinking GCD register as carry scratch (gate-hosting
    \S\ref{sec:gatehost} in its first form, reaching the published ``$9n$'' peak floor),
    square-recompute eviction (\S\ref{sec:liverange}), and \textbf{joint-pin plateau breaking}
    (\S7.9) --- the realization, present from the very start, that a peak co-owned by several scratch
    clusters drops only when you reduce them all \emph{together}.
  \item \textbf{2002 $\to$ 1698 (the engine swap):} the first true \emph{architectural} jump replaced
    the entire Kaliski stack with a \textbf{dialog-GCD inverter that streams a compressed transcript}
    (\S\ref{sec:dialog}) instead of holding wide history registers resident --- one change worth
    $-304$ qubits, and the lineage everything after descends from.
\end{itemize}

\subsubsection{Era 2 --- The dialog-GCD frontier: 1698q $\to$ 1168q (and the wall)}

On the dialog-GCD engine the peak fell steadily via the \textbf{windowed/chunked apply} dial
(\S\ref{sec:windowed}, the era's biggest peak lever), carry self-hosting on the operand lanes
(\S\ref{sec:borrowcarry}/\S\ref{sec:gatehost}), shift-free modular doublings (\S\ref{sec:shiftfree}),
keep-live-vs-recompute in the Solinas fold (\S\ref{sec:liverange}), the first \textbf{recompute-to-free
fold carries}, and progressively tighter transcript codecs:

\begin{center}
\begin{tabular}{p{3.8cm}p{4.5cm}p{5.2cm}}
\toprule
\textbf{Transition} & \textbf{Technique} & \textbf{What was released} \\
\midrule
1698$\to$1211 & windowed apply + recompute-to-free + carry self-hosting & the whole 1698$\to$1211 descent (\S\ref{sec:windowed}/\S\ref{sec:borrowcarry}/\S\ref{sec:shiftfree}) \\
1211$\to$1193 ($-18$q) & Live-range hole (\S\ref{sec:liverange}) & GCD-apply scratch block freed during shift sub-phase; reacquired after \\
1193$\to$1192 ($-1$q) & Truncation (\S7.4) & Square segment one notch tighter \\
1192$\to$1185 ($-7$q) & Transcript codec (\S\ref{sec:transcript}) & Per-step fold carry scheduling + exact 11-bit head codec \\
1185$\to$1170 ($-15$q) & Plateau decomp (\S7.9) + 5 techniques & Tail codec, streaming apply, square rebalance, 5 co-binders \\
1170$\to$1169$\to$\textbf{1168} & MBU vent on fold carry (\S\ref{sec:mbu}/8.B) & \texttt{hmr+cz\_if} on the $-f$ fold peak owner: 0 Toffoli to vent it --- then \textbf{the wall} \\
\bottomrule
\end{tabular}
\end{center}

\textbf{The 1168q wall.} The dialog-GCD frontier reached 1168q (BitWonka, 6/15) and \emph{stuck there
for days}. Every remaining lever on that engine was a single-qubit grind against a hard floor: with
the two 256-bit $Q$ coordinates materialized as resident quantum operands across the GCD peak, the
architecture had no more room. Breaking 1168 needed a different circuit, not another knob.

\subsubsection{Era 3 --- The \emph{ludicrous} breakthrough: 1167q $\to$ 1152q}

On 6/18 \textbf{tob-joe (Trail of Bits), \texttt{af5abb1}}, broke the wall by porting trailmix's
\textbf{``ludicrous''} point-add --- \emph{a new circuit family, distinct from dialog-GCD} --- onto
the challenge. It landed at \textbf{1167q $\times$ 1{,}422{,}591} (only $-1$q below the wall on day
one), but the point wasn't the $-1$q: it opened a \textbf{fundamentally lower-headroom family} the
dialog-GCD line could not reach, and over the next days drove to 1152. Its design (the architecture
\S\ref{sec:trailmix} describes in full):

\begin{itemize}
  \item \textbf{The decisive lever --- $Q$ classical \emph{and off-peak} ($-512$q).} Hold both 256-bit
    $Q$ coordinates as \emph{classical} registers and load them into a transient quantum temp
    \textbf{only at off-peak coordinate steps}, so all 512 of those qubits never co-reside at the GCD
    peak. (Earlier circuits knew $Q$ was a constant but still paid for it as resident quantum operands
    at the peak; ludicrous is what finally kept it off --- tob-joe's note calls this ``the decisive
    product-min lever.'')
  \item \textbf{One shared inversion, two passes.} The whole add is two GCD passes sharing one
    modular-inversion primitive (\texttt{Direction::\{Inverse, Forward\}}), a fused square-subtract for
    $\lambda^2$, and a zero-Toffoli negate $ox - x2 = -(x2 - ox)$ (\S9.9).
  \item \textbf{Schrottenloher jump-GCD (jump=2, 258 steps).} Each forward step records a 3-bit dialog
    symbol $(\text{subtracted}, \text{swap}, s_2)$; the reverse replays it. The GCD registers
    \textbf{shrink and regrow along a baked width schedule}, running the adders in the freed headroom
    --- this sets the 1167q floor and is the primary lossy lever.
  \item \textbf{All-triple base-5 codec} (\S\ref{sec:transcript}): 3 jump-2 symbols (9 raw bits) fold
    into a 7-bit code (only 5 of 8 patterns reachable; $\log_2(5^3)\approx 6.97$), from a
    SAT-synthesized in-place pairs core --- tighter than the earlier K5 15$\to$12 codec.
  \item \textbf{Inline interleaving + streamed decompression:} each window is compressed the instant
    its symbols are produced, so the full raw tape never exists (resident \textbf{603q vs $\sim$775
    raw}); the reverse decompresses one window at a time, and the apply is fused \emph{into} the GCD
    passes, not a separate phase.
  \item \textbf{A vented-adder zoo on baked per-call schedules} (\S9.2): plain /
    variable-chunk / headroom-adaptive / chunked-Cuccaro / carry-out adders plus a fused
    \texttt{double+cdouble} $(e+2d)\cdot f$ fold; for each of the $\sim$258$\times$2 calls the adder
    choice, carry cap, vent count and chunk split are replayed from baked tables sized for that call's
    live-qubit position.
  \item \textbf{Deliberate exact-modular truncations:} the $+f$ fold touches only the low 54 bits
    ($\sim$$2^{-21}$ miss) and the swap comparators recompute on a narrow top window --- bounded rare
    errors for large savings (\S9.18, \S\ref{sec:density}). Everything is baked, so \texttt{build()} is
    deterministic and a tail nonce (28565) is ground to land all 9024 shots clean.
\end{itemize}

From that base the trailmix\_ludicrous line ground down to the SOTA:

\begin{center}
\begin{tabular}{p{3.8cm}p{4.5cm}p{5.2cm}}
\toprule
\textbf{Transition} & \textbf{Technique} & \textbf{What was released} \\
\midrule
1167$\to$1166 & Step-0 swap\_flag elimination & \texttt{swap\_flag} becomes \texttt{Option<QubitId>}, lazy-alloc; freed one tape lane \\
1166$\to$1165 & Odd-u0 parking (\S\ref{sec:liverange}) & \texttt{u[0]} is provably 1; park+loan the lane, restore in 2 gates \\
1164$\to$1163 & Source-as-carry suffix & Paid drop: surrendered apply-phase vents to free one carry lane \\
Toffoli: $-22.4$M & Karatsuba square (\S9.6) & 3 sub-squares instead of 4 \\
Toffoli: $-1.33$M & GAP\_J2 narrowing (\S9.11) & 22-line per-step comparator table edit \\
Toffoli: $-377$k & Constprop deeper (\S9.16) & \texttt{CONSTPROP\_MAX\_ITERS} $16\to256$ \\
1163$\to$1159$\to$1156 & Headroom clamp returns (\S7.10 + shelved-lever \S\ref{sec:exchange}) & Same clamp that lost early now wins on the cheap Karatsuba$+$dead-CCX base ($\sim$527 T/qubit) \\
1156$\to$1153 & Clamp $+$ zero-cin fold $+$ FFG cap (\S7.10), paid by cinc (\S\ref{sec:cinc}) $+$ iterated dead-CCX (\S\ref{sec:deadccx}) & Qubit drop alone cost $+11{,}369$ avgT (rejected); Toffoli levers stacked back to win \\
\textbf{1153$\to$1152} & Free-and-recompute cy0 (\S7.5) & \texttt{cy0 = ctrl \& !a0}; loan lane for 40 binding folds \\
Structural & 6dafa07 pivot (\S9.8/\S\ref{sec:density}) & Replaced empirical dead-CCX with $\approx15$ structural predicates \\
\bottomrule
\end{tabular}
\end{center}

\textbf{Key meta-lessons}:
\begin{enumerate}
  \item \textbf{Every ``floor'' fell to a better encoding or a better question.} Not always a new
    algorithm --- often a denser codec or the question ``whose value is a cheap function of live
    qubits?''

  \item \textbf{The shelved-lever rule is real.} The 1157q clamp tried early and lost; won after Karatsuba.

  \item \textbf{Plateau mode is the endgame.} The final qubit drops each required coordinated
    packages of 4--9 independent local reductions applied simultaneously.

  \item \textbf{Structural correctness wins over island-exact tricks.} The 6dafa07 pivot to
    structural dead-CCX eliminated the empirical approach entirely --- correct for all inputs,
    not just the 9024.
\end{enumerate}

\subsection{The low-qubit track: the Shrunken-PZ inversion (\S\ref{sec:twocomps}, track 2)}

\subsubsection{What the Shrunken-PZ approach is (the backbone of this track)}

Almost every record on the low-qubit track is built on one architecture, \textbf{Shrunken-PZ}
(Proos--Zalka). It is worth understanding as a whole, because it is a \emph{fundamentally different
inversion engine} from the ludicrous/dialog-GCD circuit that holds the product SOTA --- different
module, different objective, sharing only the problem statement ($P \mathrel{+}= Q$).

\textbf{The core idea.} Where the ludicrous track records a small \textbf{decision tape} and replays
it (\S\ref{sec:dialog} dialog transcript), Shrunken-PZ runs the \textbf{Proos--Zalka binary-GCD /
Kaliski divstep inversion directly as a bit-by-bit reversible state machine} --- no big transcript.
It carries the actual Euclidean state in \textbf{five working registers} and \emph{resizes them every
step}:

\begin{center}
\begin{tabular}{p{2.0cm}p{5.0cm}p{6.0cm}}
\toprule
\textbf{register} & \textbf{meaning} & \textbf{width over the 530 steps} \\
\midrule
$A$, $B$ & the two running GCD magnitudes & start 256, \textbf{shrink toward 1} as the GCD converges \\
$CA$, $CB$ & the two B\'ezout cofactors (inverse accumulators) & start 1, \textbf{grow toward 256} \\
$Q$ & the per-step quotient & small, $\sim$21--25 bits \\
\bottomrule
\end{tabular}
\end{center}

\textbf{Why the peak is in the middle, and why register sharing fits.} $A,B$ shrink while $CA,CB$
grow, so the qubit peak is not at either end --- it is at the \textbf{GCD crossover} ($\sim$step 363),
where both pairs are mid-sized:
\[
  \text{peak} \approx 2\max(|A|,|B|) + 2\max(|CA|,|CB|) + |Q| + \text{FIXED}(\sim267).
\]
At the crossover this is $[81,81,248,248,23] = 681$ working bits $+ \sim$267 fixed (the 256-bit
coordinate passengers, modulus, sign/parity, counter, comparator scratch) $=$ \textbf{948 qubits}.
Crucially $A$ and its cofactor $CA$ are \textbf{complementary in size} (one shrinks exactly as the
other grows) --- which is what makes EEA register sharing (\S\ref{sec:regshare}) able to pack
$A\Vert CA$ into one word and push toward 851/829.

\textbf{The levers that define the track} (each is a \S7 technique, applied to the divstep machine):
\begin{itemize}
  \item \textbf{Dynamic-width registers (\S7.6)} --- resize $A,B,CA,CB$ per step; free known-zero
    high bits.
  \item \textbf{CLZ-context window narrowing} --- the divstep arithmetic only touches bits above a
    provably-zero low region, so each op runs over a narrowed window $[lo..len)$. The ``robust width
    envelope'' projects the tightest per-step window that still covers all sampled inputs --- the
    qubit-side twin of empirical dead-CCX (\S\ref{sec:deadccx}), and equally \textbf{island-exact}.
  \item \textbf{Known-constant teardown (\S7.7)} --- at convergence $A{=}0, B{=}1, CA{=}p, Q{=}0$ are
    constants; XOR them out (0 Toffoli) before the $\lambda$ multiply so they don't co-reside with it.
  \item \textbf{Passenger ghosting (\S7.14)} --- HMR-ghost $dy$/$\lambda$ so only one 256-bit
    coordinate co-resides with the EEA peak ($+256$, not $+512$).
  \item \textbf{Pipelined two-quotient divstep} --- build the \emph{next} quotient while draining the
    \emph{previous}, so the quotient record stays $\sim$26 bits instead of a full 256-bit tape.
  \item \textbf{EEA register sharing (\S\ref{sec:regshare})} --- the deepest lever; packs operand and
    cofactor, reaching 851/829q.
\end{itemize}

\textbf{Why it is the qubit backbone but not score-competitive.} Running full per-step modular
arithmetic \emph{without} the transcript-replay and Karatsuba machinery makes the divstep inversion
inherently gate-heavy: $\sim$33--55M Toffoli at 948--980q ($\sim$40$\times$ the ludicrous SOTA's
1.36M), and $\sim$460M at 851q. So Shrunken-PZ maps the \textbf{qubit floor} (objective 2) but loses
the \textbf{product} by a wide margin --- a qubit-lower-bound witness, not a contender on score.

\subsubsection{The record ladder}

Running \emph{alongside} the Q$\times$T product-SOTA journey above is this separate \textbf{low-qubit
competition} line --- the Shrunken-PZ family --- whose objective is to minimize peak qubits with
Toffoli unconstrained:

\begin{center}
\begin{tabular}{p{2.6cm}p{6.0cm}p{4.4cm}}
\toprule
\textbf{Qubit record} & \textbf{Technique} & \textbf{Notes} \\
\midrule
1050q & trailmix shrunken-PZ embedded op stream & early low-qubit route \\
1019q & source-level dynamic-width PZ port & \\
988$\to$956q & dirty-borrow / gate-hosting lever stack & \\
952$\to$948q & further lane borrowing + thin schedule & lowest \emph{fully-validated} witness \\
\textbf{851q} (\texttt{e7dd3de}, 6/23) & \textbf{EEA register sharing (\S\ref{sec:regshare}), Luo 2026} & \textbf{$-97$q; the breakthrough} \\
829q (\texttt{1dd61ca}, 6/24) & more borrowing/relocation on a frozen envelope & current low-qubit frontier \\
\bottomrule
\end{tabular}
\end{center}

These are records on the \textbf{pure-qubit objective}, where the $\sim$400--500M Toffoli they spend
simply does not enter the scoring. The 851q and 829q figures come from an analysis-oracle accounting
(\S\ref{sec:regshare}) --- read them as the leading \textbf{lower-bound witnesses} for the qubit
floor; the 948q route is the lowest fully-validated circuit on the track. This is a \emph{different
game} from the Q$\times$T product SOTA (1152q $\times$ 1.36M, value-exact): more qubits but
$\sim$340$\times$ fewer Toffoli. The two minima are reached by different constructions, and which one
is ``best'' depends entirely on which competition you are entering --- see \S\ref{sec:regshare} and
\S\ref{sec:exchange}.

\subsection{The Pareto-frontier track: clean (Q, T) basis circuits 1153q $\to$ 1133q (\S\ref{sec:twocomps}, track 3)}
\label{sec:paretopush}

The third track (\S\ref{sec:twocomps}, objective 3) maps the \textbf{(Q, T) Pareto frontier} below the
1153q SOTA as a sequence of \textbf{clean, dead-CCX-free basis circuits} others can fork.
``Dead-CCX-free'' is the precise claim: every rung hard-sets \texttt{DROP\_DEAD\_ROBUST\_DISABLE=1},
omitting the most overfit lever (whose hard-coded gate-index \texttt{.idx} lists are tied to one exact
op stream, \S\ref{sec:density}) --- that omission is what keeps the basis \emph{reusable}. They are
\textbf{not fully value-exact}, however: each still carries the route's standard calibrated
\textbf{approximations} (comparator-width narrowing \S9.18, carry truncation \S7.4, and at the lowest
rungs \textbf{active-width trimming} \S\ref{sec:density}) that are island-exact and require a per-rung
nonce hunt. So treat the curve as a \emph{cleaner, more forkable} reference than a dead-CCX-stacked
SOTA, not as an all-inputs-provable bound.

Unlike the low-qubit witness track above, these stay in the \textbf{useful corner} --- every point
beats the published \emph{single point-addition} estimate on \emph{both} axes: $q < 1175$ and
$T \approx$ 1.37--1.46M, comfortably under the $\sim$2.69M ($2^{21.36}$) Toffoli of the Babbush
space-optimized secp256k1 point-add (Schrottenloher 2026, Table 1). So the whole curve clears the bar
with $\sim$20--40 qubits and roughly 2$\times$ Toffoli to spare. They are ``rejected'' on the
Q$\times$T product \emph{by construction} (they trade qubits down for Toffoli up); the deliverable is
the exchange \emph{curve}, not a score.

\begin{center}
\begin{tabular}{p{1.1cm}p{2.2cm}p{2.6cm}p{6.0cm}}
\toprule
\textbf{q} & \textbf{$T_{\text{clean}}$} & \textbf{exch.\ vs prev} & \textbf{new lever introduced} \\
\midrule
1153 & 1,392,603 & --- (anchor) & dead-CCX off baseline \\
1147 & 1,396,769 & $-694$ T/q & fold-vent clamp (\texttt{TLM\_FOLD\_CALL\_CODE\_OVERRIDES}) \\
1146 & 1,408,582 & $-11{,}813$ T/q & graduated$\to$dirty carry-in (\texttt{TLM\_GRAD\_DISABLE}) \\
1143 & 1,411,643 & $-1{,}020$ T/q & codec scratch reuse $+$ direct-dirty fold ctl \\
1142 & 1,415,790 & $-4{,}147$ T/q & pair-raw codec last-k $+$ tighter layout search \\
1141 & 1,423,723 & $-7{,}933$ T/q & no-ancilla/dirty body $+$ constant-lane loan family \\
\textbf{1133} & \textbf{1,460,511} & $-4{,}599$ T/q & \textbf{GCD active-width trimming} $+$ fold-cout drop \\
\bottomrule
\end{tabular}
\end{center}

The qubits sort by exchange rate into four lever classes: \textbf{fold-vent clamping} (the cheap
qubits --- clamp a peak fold's measured-vent windows, \S9.2), \textbf{no-ancilla / dirty-borrow
substitution} (convert a carry ancilla into Toffoli, cf.\ \S\ref{sec:borrowcarry}),
\textbf{constant-lane loan $+$ recompute} (the Bennett pebble move on provably-constant GCD low bits,
\S7.1/\S\ref{sec:pebbling}), and --- new at 1133 --- \textbf{GCD active-width trimming}: per-step
shrinking of the live GCD operand width in the convergence tail
(\texttt{TLM\_GCD\_ACTIVE\_WIDTH\_TRIM=3} after step 205), the dynamic-width-register idea (\S7.6)
applied to the ludicrous GCD. The 1133 rung is a \emph{value-exact repair} of an over-aggressive 1131
attempt: trimming the active width is a \textbf{graded} correctness lever (\S\ref{sec:qubitloop}) ---
push until the residual breaks, then give back one notch (1131$\to$1133) so a clean nonce is
reachable. The push also introduced a reusable general pass --- \textbf{reset-bounded qubit-id
compaction} (\S\ref{sec:compact}). Full analysis:
\href{https://github.com/jieyilong/ecdsafail_skills/blob/main/ecdsafail-circuit-optimization/references/pareto-frontier-push.md}{\texttt{pareto-frontier-push.md}}.

% ===========================================================================
\section{Open Frontiers}
\label{sec:frontiers}
% ===========================================================================

\subsection{Density-neutral Toffoli cuts not yet applied}

\begin{enumerate}
  \item \textbf{Gidney 2025 streaming vented adder} (arXiv:2507.23079): Uses 2 clean $+$ $(n-2)$
    dirty ancilla for $\approx3n$ Toffoli. Partially implemented in \texttt{venting.rs} but
    \textbf{leaks phase} --- dirty ancilla phase corrections are incomplete. Expected $\approx n/4$
    Toffoli savings per adder once fixed.

  \item \textbf{Apply-swap structural dead-CCX}: The GCD apply-step \texttt{cswap} block
    ($\approx256$ cswap per iter $\times$ 258 iters $\times$ 2 passes) has NO per-bit
    structural-dead skip. Porting the structural-dead analysis to these operations is a
    high-priority open lever.

  \item \textbf{Extended Cuccaro structural dead-carry}: Currently checked for call indices 13--37
    only. Extending to ALL Cuccaro adder call sites would save proportionally more.

  \item \textbf{Apply-swap involutory pair cancellation}: The forward GCD-swap and the immediately
    following GCD-cswap on the same $(u/v)$ lanes may cancel algebraically (swap twice $=$ identity).
    Not yet checked structurally.

  \item \textbf{Windowed arithmetic / QROM lookup} (Gidney 2019, arXiv:1905.07682) --- \emph{the
    biggest untried structural lever.} A multiply-by-a-quantum-value normally does $2^w$ controlled
    add-shifts. The windowed trick reads a \textbf{$w$-bit window} of the multiplier and uses it to
    \textbf{address a QROM table lookup} that adds the precomputed partial product in one shot ---
    turning $2^w$ controlled adds into one table read. The B\'ezout-reconstruction multiplies
    (\S\ref{sec:dialog}) are $\sim$90\% of the Toffoli, so a windowed/QROM form there is the largest
    single cut on the board. Open question: whether the value-dependent reconstruction multiply can
    be windowed the way the fixed-base $k\cdot G$ scalar mult already is.
\end{enumerate}

\subsection{Density-neutral qubit cuts not yet applied}

\begin{enumerate}
  \item \textbf{The co-bind plateau at 1152q}: Multiple independent phases all tie at 1152q.
    Each needs a local reduction before the global peak drops.

  \item \textbf{FFG cy0-style recompute on other binders}: The cy0 trick (\S7.5) worked because
    \texttt{cy0} was a cheap function of live qubits. Other binder qubits may have similar
    cheap-function properties not yet identified.

  \item \textbf{SumHiLo/shifted vents} (\texttt{TLM\_SQUARE\_SUMHILO\_VENT},
    \texttt{TLM\_SQUARE\_VENT\_SHIFTED}): Exist in the codebase as default-OFF options. Off-peak
    phases could enable these for free Toffoli savings.
\end{enumerate}

% ===========================================================================
\section{Quick-Reference Summary Tables}
\label{sec:tables}
% ===========================================================================

\subsection{Qubit reduction techniques}

\begin{center}
\small
\begin{tabular}{p{4.8cm}p{3cm}p{2.8cm}p{2.5cm}}
\toprule
\textbf{Technique} & \textbf{Qubits saved} & \textbf{Toffoli cost} & \textbf{Density-neutral} \\
\midrule
Live-range hole (\S7.1) & $k$ qubits & $+2 \times C_V$ & Yes \\
Passenger relocation (\S7.2) & $k$ qubits & 0 & Yes \\
Range-bound guard removal (\S7.3) & 1 per register & 0 & Yes \\
Truncation (\S7.4) & varies & 0 & Partial \\
Free-and-recompute cheap value (\S7.5) & 1 & $\approx0$ (CNOTs) & Yes \\
Dynamic-width registers (\S7.6) & varies & small & Partial \\
Known-constant teardown (\S7.7) & $k$ & 0 (CNOT $=$ Clifford) & Yes \\
Transcript compression (\S7.8) & varies & decode overhead & Yes \\
Plateau decomposition (\S7.9) & 1 (global) & combined cost & Yes \\
Dynamic headroom clamp (\S7.10) & 1 per ceiling drop & small & Yes \\
Lazy carry liveness (\S7.11) & 1/chunk & 0 & Yes \\
In-place decode (\S7.12) & varies & codec overhead & Yes \\
Borrowed-carry fusion (\S7.13) & $k$ & 0 & Yes \\
HMR ghosting (\S7.14) & 256 & $+1$ modular multiply & Yes \\
EEA register sharing (\S7.15) & $\sim$97--119 (to 851$\to$829q) & high (irrelevant to low-q objective) & Yes --- low-qubit competition \\
Gate-hosting (\S7.16) & 1 per hosted gate & 0 (clean) / $+1$ Tof (dirty) & Yes if host-zero \emph{proven}; island-exact if \emph{sampled} \\
Reset-bounded id compaction (\S7.17) & gap (max-id $-$ true peak) & 0 & Yes (relabels wires only) \\
Windowed/chunked materialization (\S7.18) & many (the \texttt{F\_CUT} dial) & $+$few (boundary cmps) & Yes \\
Shift-free modular doublings (\S7.19) & $\sim$24 (drops spill$+$flags) & $+$few & Yes \\
Classical constants off-peak (\S7.20) & \textbf{512} (the Q-classical lever) & 0 & Yes \\
GCD active-width trim (\S7.6 applied) & varies (deep tail) & small & Partial (graded; nonce-hunted) \\
\bottomrule
\end{tabular}
\end{center}

\subsection{Toffoli reduction techniques}

\begin{center}
\small
\begin{tabular}{p{4.8cm}p{3cm}p{2.5cm}p{2.5cm}}
\toprule
\textbf{Technique} & \textbf{Toffoli saved} & \textbf{Qubit cost} & \textbf{Density-neutral} \\
\midrule
MBU / carry vent (\S9.1/2) & 1 per vent & $+1$ temporary & Yes \\
Conditional replay (\S9.3) & $(1-f)\times\text{block}$ & 0 & Yes \\
Algebraic fusion (\S9.4) & varies & 0 & Yes \\
Witness-first dataflow (\S9.5) & $\approx13$M / instance & 0 & Yes \\
Karatsuba square (\S9.6) & $\approx22$M (one change) & 0 & Yes \\
NAF recoding (\S9.7) & $\approx5$--10\% of fold cost & 0 & Yes \\
Pseudo-Mersenne reduction (\S9.18) & huge (15\% vs 34\% of CCX) & 0 & Yes (fold); Partial (lsbs/msbs trunc) \\
Merge adjacent controlled-swaps (\S9.19) & $n$ per merged pair ($-274$k once) & 0 & Yes \\
Structural dead-gate skip (\S9.8) & varies & 0 & Yes \\
Classical constant ops (\S9.9) & significant & 0 & Yes \\
Converged-tail elision (\S9.10) & hundreds/iter $\times K$ & 0 & No \\
GAP\_J2 narrowing (\S9.11) & $\approx1.33$M / 22-line edit & 0 & No \\
Ancilla-free majority (\S9.12) & 1 per majority & 0 & Yes \\
Exact-adder recovery (\S9.13) & $\approx2{,}600$/call & 0 & Yes \\
Redundant cleanup skip (\S9.14) & varies & 0 & Risky \\
Off-peak loosening (\S9.15) & free refund & 0 & Yes \\
Constprop deeper (\S9.16) & $\approx377$k & 0 & Yes \\
Cinc cascade replacement (\S9.17) & $\approx5{,}175$ ($O(t^2)\to O(t)$) & 0 & Yes \\
Empirical dead-CCX drop (\S9.8) & $\approx13{,}000+$ (1st pass) & 0 & No (island-overfit) \\
Iterated 2-pass dead-CCX (\S9.8) & $\approx2{,}400$/extra pass & 0 & No (island-overfit) \\
\bottomrule
\end{tabular}
\end{center}

\subsection{Key theoretical results}

\begin{center}
\small
\begin{tabular}{p{4.5cm}p{4.5cm}p{4.5cm}}
\toprule
\textbf{Result} & \textbf{Source} & \textbf{Insight} \\
\midrule
MBU: uncompute AND for 0 $T$-gates & Jones 2013 (Phys Rev A); Gidney 2018 (arXiv:1709.06648) & $X$-measurement $+$ CZ $=$ free uncompute \\
1 Toffoli $=$ exactly 7 $T$-gates & Gosset et al.\ 2013 (arXiv:1308.4134) & Tight lower bound \\
Bennett classical pebbling & Bennett 1989 (SIAM J.\ Comput.\ 18:4) & Classical: exponential constant in space savings \\
Quantum spooky pebbling & Kornerup, Sadun, Soloveichik 2021 (arXiv:2110.08973) & Quantum: polynomial constant (exponential advantage) \\
Parallel spooky pebbling & Kahanamoku-Meyer et al.\ 2025 (arXiv:2510.08432) & $2.47 \times \log(\ell)$ qubits for $\ell$-step chain \\
Conditionally-clean ancilla & Khattar \& Gidney 2024 (arXiv:2407.17966) & $3n$ Toffoli MCX with $\log^* n$ clean ancilla \\
secp256k1 \emph{full ECDLP} frontier & Babbush, Gidney et al.\ 2026 (arXiv:2603.28846) & $\leq1200$q / $\leq90$M Toffoli for the whole Shor run ($\sim$28 windowed PAs) \\
secp256k1 \emph{single point-add} (the challenge unit) & Babbush space-opt, via Schrottenloher 2026 Table 1 (arXiv:2606.02235) & \textbf{1175q / $2^{21.36}\approx2.69$M Toffoli} (the ``useful corner''); Schrottenloher space-opt $=$ 1192q / $2^{21.19}\approx2.39$M \\
Karatsuba algorithm & Karatsuba \& Ofman 1962 & $O(n^{1.585})$ vs $O(n^2)$ for multiplication \\
EEA register sharing (origin) & Proos \& Zalka 2003 (arXiv:quant-ph/0301141) & Inversion dominates point-add space; operand$+$cofactor packing \\
Compact exact reversible inversion & Luo et al.\ 2026 (arXiv:2604.02311) & $3n + 4\lfloor\log_2 n\rfloor$ qubits (1333q, $n{=}256$) via bit-length-tracked sharing \\
DQI ``Dialog'' in-place EEA multiply (\S\ref{sec:dialog}) & Khattar, Shutty \& Gidney 2025 (arXiv:2510.10967) & construct/reconstruct split; $(y,0)$ seeding fuses inverse$+$multiply, no explicit inverse \\
Windowed arithmetic / QROM (\S\ref{sec:frontiers}) & Gidney 2019 (arXiv:1905.07682) & $2^w$ controlled adds $\to$ one $w$-bit-window table lookup \\
\bottomrule
\end{tabular}
\end{center}

\vfill
\noindent\rule{\textwidth}{0.4pt}
\vspace{2pt}

\noindent\textit{See also} --- these detailed reference docs live in the GitHub repo
\href{https://github.com/jieyilong/ecdsafail_skills}{\texttt{github.com/jieyilong/ecdsafail\_skills}},
under \texttt{ecdsafail-circuit-optimization/references/}:
\begin{itemize}\small
  \item \href{https://github.com/jieyilong/ecdsafail_skills/blob/main/ecdsafail-circuit-optimization/references/density_neutral_tradeoffs.md}{\texttt{density\_neutral\_tradeoffs.md}} --- detailed technique catalog with exchange rate math
  \item \href{https://github.com/jieyilong/ecdsafail_skills/blob/main/ecdsafail-circuit-optimization/references/gidney-techniques.md}{\texttt{gidney-techniques.md}} --- Gidney's full lever catalog with blog/paper pointers
  \item \href{https://github.com/jieyilong/ecdsafail_skills/blob/main/ecdsafail-circuit-optimization/references/external-literature-2000-2026.md}{\texttt{external-literature-2000-2026.md}} --- broader academic literature map
  \item \href{https://github.com/jieyilong/ecdsafail_skills/blob/main/ecdsafail-circuit-optimization/references/REPORT_1168_wall_revamp.md}{\texttt{REPORT\_1168\_wall\_revamp.md}} --- the trailmix\_ludicrous introduction
  \item \href{https://github.com/jieyilong/ecdsafail_skills/blob/main/ecdsafail-circuit-optimization/references/frontier-1211-to-1170.md}{\texttt{frontier-1211-to-1170.md}} --- detailed 1211$\to$1170 step-by-step record
  \item \href{https://github.com/jieyilong/ecdsafail_skills/blob/main/ecdsafail-circuit-optimization/references/pareto-frontier-push.md}{\texttt{pareto-frontier-push.md}} --- the (Q,T) Pareto-frontier push (\S\ref{sec:paretopush})
  \item \href{https://github.com/jieyilong/ecdsafail_skills/blob/main/ecdsafail-circuit-optimization/references/SHRUNKEN_PZ_q948_track.md}{\texttt{SHRUNKEN\_PZ\_q948\_track.md}} --- the low-qubit Shrunken-PZ track (\S\ref{sec:twocomps})
\end{itemize}

\vspace{4pt}
\noindent\textbf{Note on code paths.} Source files named in this primer (e.g.\ \texttt{arith.rs},
\texttt{fused.rs}, \texttt{register\_shared\_eea.rs}) refer to the \textbf{ecdsa.fail challenge
circuit} --- the \texttt{trailmix\_ludicrous} / \texttt{trailmix\_port} modules --- pulled
per-submission via the \texttt{ecdsafail} CLI (\texttt{ecdsafail reset <submission>}); they are
\emph{not} part of this docs repo.

\end{document}
